% From Scrap to Value -- the document itself. Do not compile this file: compile one of its two wrappers.
%   manuscrit.tex        reads genere/  (demonstration data; the folder and its PDF are versioned)
%   manuscrit_site.tex   reads site/    (your data; never versioned, never packaged)
% Both folders are written by:  python -m erc.papier   (nothing in them is typed by hand; see LISEZ-MOI.md)
\documentclass[11pt]{article}
\usepackage[T1]{fontenc}
\usepackage[utf8]{inputenc}
\usepackage{mathptmx}
\usepackage[margin=1in]{geometry}
\usepackage{amsmath,amssymb,amsthm,mathtools}
\usepackage{booktabs,tabularx,array}
\usepackage{graphicx}
\usepackage{float}
\usepackage[font=small,labelfont=bf]{caption}
\usepackage{setspace}
\usepackage{xcolor}
\usepackage[round,authoryear]{natbib}
\usepackage[colorlinks=true,linkcolor=black,citecolor=black,urlcolor=black]{hyperref}

\input{\genere/macros}

% ---- statements
\newtheorem{proposition}{Proposition}
\newtheorem{lemma}{Lemma}
\newtheorem{corollary}{Corollary}
\theoremstyle{definition}
\newtheorem{assumption}{Assumption}
\newtheorem{definition}{Definition}
\theoremstyle{remark}
\newtheorem{remark}{Remark}

% ---- notation
\newcommand{\cI}{\mathcal{I}}
\newcommand{\cJ}{\mathcal{J}}
\newcommand{\cD}{\mathcal{D}}
\newcommand{\cL}{\mathcal{L}}
\newcommand{\cS}{\mathcal{S}}
\newcommand{\cC}{\mathcal{C}}
\newcommand{\bC}{\bar{\mathcal{C}}}
\newcommand{\E}{\mathbb{E}}
\renewcommand{\Pr}{\mathbb{P}}
\newcommand{\R}{\mathsf{R}}
\newcommand{\Sd}{\mathsf{S}}
\newcommand{\out}{\varnothing}
\newcommand{\pos}[1]{\left(#1\right)^{+}}
\DeclareMathOperator{\EVPI}{EVPI}
\DeclareMathOperator{\Bfun}{B}

% ---- body of a generated table (the primitive \input is expandable, which a tabular needs before \bottomrule)
\makeatletter
\newcommand{\corps}[1]{\@@input \genere/#1.tex }
\makeatother

% ---- what remains to be written once the site data have been run
\newcommand{\apres}[1]{\textcolor{red!60!black}{\textbf{[To write after the run on site data:} #1\textbf{]}}}
\newcommand{\bandeau}{\ifnum\DonneesDemo=1\relax
  \begin{center}\fbox{\parbox{0.92\textwidth}{\small\textbf{Demonstration data.} Every case-study number in this
  draft (Section~\ref{sec:case}, and the sentences of the abstract and introduction that quote it) was produced from
  \emph{fictional} data shipped with the software. They show what the system computes; they are not results. Running
  \texttt{python -m erc.papier} on the site's data replaces them all and removes this box.}}\end{center}\fi}
\newcommand{\noteechelle}{\ifnum\EchelleMasquee=1\relax Monetary amounts of the case study are multiplied by an
  undisclosed constant. \fi}

\title{\textbf{From Scrap to Value: Pricing Labour, Learning Yields, and Buying Knowledge\\ for the Disposition of
Scrapped Service Units}}
\author{[Author names withheld in this draft]\\ {\small [Affiliation]}}
\date{Draft of \today}

\begin{document}
\maketitle
\bandeau

\begin{abstract}
\noindent A repair centre for medical equipment scraps the field-replaceable units it cannot repair. They leave whole,
through a bulk waste contract, although they contain components that could be reused in repairs or sorted for
specialised recyclers. Deciding what to do with them is hard for a reason that optimisation alone does not address:
most of the facts the decision needs are missing. The share of components that are still good is unknown, the flow of
scrapped units is irregular, and the material composition of components, which sets recycling prices, is recorded
nowhere. We present a decision-support system that treats the three difficulties together. Disposition is a linear
multiple-choice knapsack over unit types; we show that one number, the value of an operator hour, prices every unit
without enumerating its plans, yields a bench-level threshold rule for reuse, and bounds, from data alone, what
buffering, forecasting or reinforcement learning could add. Unknown yields are learned by posterior sampling. Missing
compositions are acquired by a chain of sources of learned reliability (rules, analogy, documents, and a language
model hosted on the owner's machine) in which every paid step must be covered by an upper bound on what the answer
can be worth; we prove the bound, the budget invariant, and the consistency of the fusion rule. Agents running on the
same self-hosted open model drive the system end to end and write a decision note whose numbers are checked against
computed results. In the case study, \CasTypes{} unit types are modelled; the system estimates an annual result of
\CasResultat{} against the current contract, a value of an additional operator hour of \CasMu, and a cost of flow
variability of \CasGammaRel{} of the fluid value. \noteechelle
\end{abstract}

\noindent\textbf{Keywords:} product recovery; disposition; multiple-choice knapsack; Lagrangian duality; value of
information; Thompson sampling; fluid bound; large language models; agents; circular economy; service parts.

\onehalfspacing

\section{Introduction}\label{sec:intro}

A repair centre for medical imaging equipment receives field-replaceable units from hospitals: power supplies,
amplifiers, control boards, detector modules. Most are repaired and returned to service. The others are found beyond
economic repair and are \emph{scrapped}. At the site that motivated this work, scrapped units leave whole, through a
bulk waste contract that produces no revenue. Yet a scrapped unit is rarely a worthless one. It failed for one cause,
and the components that did not cause the failure are often sound: they can be extracted, tested and \emph{reused}
in later repairs, where they replace parts that are costly and slow to procure. What cannot be reused can be
\emph{sorted} by material and sold by the kilogram to specialised recyclers, who pay for printed circuit boards,
copper or aluminium many times what a mixed stream is worth. The alternative to the bulk contract is therefore a small
teardown line: open the unit, take out what is worth taking, test it, sort the rest.

Whether to install such a line, and how to run it, looks like a textbook disposition problem. It is not, and the
reason is the subject of this paper. A disposition model needs, for every component of every unit type, the
probability that it is still good, the time to extract and test it, its value if reused, its mass, and the price a
recycler pays for its material class. At the start of the project almost none of these were known.
\begin{itemize}
\item \emph{Yields are unknown.} Nobody has torn down scrapped units, so nobody knows which share of their components
  would pass a test. Consumption records of repairs give a prior, biased in directions that can be named but not
  quantified without teardown data.
\item \emph{Arrivals are irregular.} Scrapped units arrive a few per week per type, with many empty weeks, while
  operator hours are booked in advance.
\item \emph{Compositions are missing.} The bill of materials gives a company part number and a free-text description.
  The material class of a component, which fixes its recycling price and its regulatory route, is recorded nowhere.
  It can be found: by a rule on the description, by analogy with a known part, in a manufacturer's material
  declaration, by reading a data sheet. Each of these has a cost and can be wrong, and there are thousands of parts.
\item \emph{The data are private.} Part numbers and descriptions of in-house designs cannot leave the company, which
  excludes sending them to a hosted language model, the tool that would otherwise read descriptions and data sheets
  at scale.
\end{itemize}

The system we describe treats these difficulties as one problem with three layers. The \emph{disposition} layer
decides, for each unit type, whether to open it and what to do with each component, under a labour budget. The
\emph{learning} layer estimates yields and arrivals while decisions are being taken. The \emph{knowledge} layer
decides which missing facts are worth acquiring, by which means, and when to stop. The layers are coupled by a single
quantity. The dual price of labour computed by the disposition layer prices every unit and every component; the same
optimisation model, evaluated with and without a fact, gives the value of knowing it; and that value is the budget
the knowledge layer may spend to acquire it. The result is a system in which a language model is called on a
component only when the answer can pay for the call, and a human is asked only where the stake exceeds the cost of
asking.

\paragraph{Contributions.} The paper makes four contributions.

\emph{(1) Structure of the disposition problem.} The problem is the linear relaxation of a multiple-choice knapsack
over unit types, with one item per plan (a disposition for each component) and a set-up that depends on how far the
unit is opened. We show that, at a given labour price, the best plan of a unit is found by pricing its components
independently for each access level (Proposition~\ref{prop:pricing}), which removes the enumeration of plans and makes
units of a thousand components as easy as units of five. The optimal labour price coordinates all decisions
(Propositions~\ref{prop:dual} and~\ref{prop:split}), and gives the operator a threshold on the observed good rate of
a component above which it is worth extracting and testing (Proposition~\ref{prop:threshold}). With irregular
arrivals, the gap between the model on the average flow and the model replayed period by period is an upper bound on
what \emph{any} policy that stores units, forecasts arrivals or is trained by reinforcement learning can gain over
the myopic policy (Proposition~\ref{prop:fluid}). The bound is computable from the history, so the question ``should
we build a dynamic policy?'' is answered by a number before anything is built.

\emph{(2) A budgeted knowledge layer.} We define the value of establishing the material class of a component as the
difference between two optimal values of the disposition model, and bound it by a quantity computable from the
bill of materials alone (Proposition~\ref{prop:class}). The acquisition chain spends on a component only within that
bound. We prove that the chain never spends on a component more than its class can be worth
(Lemma~\ref{lem:budget}), that an acquisition excluded by the screen on a component not yet paid for has negative
expected net value (Proposition~\ref{prop:screen}), that acceptance without review is the cheaper decision whenever the value at risk
is below the cost of a review (Lemma~\ref{lem:accept}), that the fusion of sources identifies the true class at an
exponential rate whenever a source's reliability differs from chance (Proposition~\ref{prop:fusion}), and that
reliabilities estimated from human validations are unbiased only if the validations are stratified by the reason the
component was validated (Proposition~\ref{prop:strata}). In this layer a language model is a source with a learned
reliability, never ground truth.

\emph{(3) Agents on a self-hosted open model.} The whole system is driven by three agents (a pilot, a writer, a
reviewer) that run on an open-weight language model hosted on the data owner's machine. They can only call a closed
list of the system's own functions; every number of the decision note they write is checked by a program against
the computed results (Lemma~\ref{lem:grounding}); without a model, a fixed plan produces the same note from a
template. The instruction given to the model for reading part descriptions is itself improved from the human
validations by reflective prompt evolution, with a call budget and a held-out test
(Lemma~\ref{lem:heldout}). No weight is trained and no record leaves the site.

\emph{(4) An open, tested implementation and its application.} All of the above is implemented in a single
package with a test suite that checks, among other things, the optimum against an exact linear-programming solver,
every proposition on random instances, and the confidentiality of outgoing queries.
Section~\ref{sec:experiments} measures the algorithms on synthetic instances, and Section~\ref{sec:case} reports
the application to the repair centre.

\paragraph{Status of the evidence.} We separate three kinds of statement. Mathematical results are proved in
Appendix~\ref{app:proofs} and checked numerically. Measurements of the algorithms use synthetic instances and say
nothing about the site. Statements about the site come from its data. \ifnum\DonneesDemo=1\relax In this draft the
case study runs on the fictional demonstration data shipped with the software, and every sentence that depends on
field results is marked as still to be written.\fi{} We make no claim that is not of one of these three kinds.

\paragraph{Organisation.} Section~\ref{sec:literature} reviews the literature. Section~\ref{sec:model} states the
disposition model and Section~\ref{sec:structure} its structure. Section~\ref{sec:learning} treats unknown yields and
irregular arrivals, Section~\ref{sec:knowledge} the acquisition of missing compositions, and Section~\ref{sec:system}
the agents and the deployment. Sections~\ref{sec:experiments} and~\ref{sec:case} give the computational study and
the case study. Section~\ref{sec:discussion} discusses managerial insights, limits and the transfer to production.

\section{Related literature}\label{sec:literature}

\paragraph{Product recovery and disposition.} Recovery options (reuse, repair, remanufacturing, cannibalisation,
recycling) are classified by \citet{thierry1995}; quantitative models are reviewed by \citet{fleischmann1997} and
\citet{guide2009}. Choosing how far to disassemble a product and which option to apply to each sub-assembly is the
recovery-strategy problem of \citet{krikke1998}, extended to several disassembly processes and to partial
disassembly by \citet{teunter2006}. \citet{meacham1999} optimise the disassembly configuration of one product and
then, by column generation, how many units of each product follow each configuration; that programme has the form
of ours. Random yields of good parts from disassembled cores are treated by \citet{inderfurth2006}; grading of
returns by condition by \citet{ferguson2009}; the choice between remanufacturing a return and dismantling it for
spare parts by \citet{ferguson2011}, in the setting described by \citet{fleischmann2003} of returned equipment as a
source of service parts, which is the closest to a repair centre. Our disposition model belongs to this family. What
we add to it is structural: the pricing of a unit by access level under a single labour price, the threshold rule
that follows, and the reading of the variability gap as a certificate. We describe these papers from their abstracts
and bibliographic records; see the note at the end of this section.

\paragraph{Knapsack structure and Lagrangian pricing.} The disposition programme is the linear multiple-choice
knapsack problem, whose solution structure (one fractional class, items on the upper hull) is classical
\citep{sinha1979,zemel1980}. Pricing the budget and solving the separable remainder is the generalised Lagrange
multiplier method of \citet{everett1963}. Proposition~\ref{prop:split} restates this structure for several coupling
constraints; we claim no novelty for it.

\paragraph{Learning yields while deciding.} \citet{reveliotis2007} proposes learning-based disassembly planning
under uncertain part quality. Our yield learning is posterior sampling \citep{thompson1933,russo2014} applied to a
linear programme, a scheme used by \citet{ferreira2018} for network revenue management; the problem is a
combinatorial semi-bandit, for which \citet{wang2018} analyse Thompson sampling. In bandits with knapsacks
\citep{badanidiyuru2018bandits} the budget is consumed over the horizon, whereas our labour budget is renewed each period.
We do not claim a regret bound: we have not verified that the assumptions of these analyses cover our observation
structure, in which a component is observed only when it is reused. A knowledge-gradient policy
\citep{frazier2008} is a natural alternative that we have not implemented.

\paragraph{Fluid bounds.} The comparison of a stochastic system with its deterministic counterpart on average data
is the upper bound of revenue management \citep{gallego1994} and rests on the Jensen inequalities of stochastic
linear programming \citep{madansky1960}. Proposition~\ref{prop:fluid} applies the argument to disposition under a
renewed budget and reads the resulting gap as a decision rule about modelling effort.

\paragraph{Value of information.} Information value theory goes back to \citet{howard1966}. In closed-loop supply
chains, the value of knowing yields before deciding is studied by \citet{ferrer2004} and \citet{ketzenberg2006}.
We compute the value of a fact through the operational model that will use the answer, and we distinguish a
parameter that can be acted on at its mean (a yield, a price) from a fact without which an option does not exist
(the class of a component), for which no ``decision at the mean'' is available.

\paragraph{Combining fallible sources.} Estimating the truth and the error rates of several observers jointly is
the problem of \citet{dawid1979}. Our fusion rule is their model with one accuracy parameter per source type and
uniform errors, with two differences: accuracies are shared by source type and learned from human-validated cases
rather than by expectation--maximisation, and they are integrated out, so that several statements from one type of
source are worth less than their product. The stratified estimation of accuracies is survey sampling
\citep{horvitz1952} applied to validations that are not a random sample.

\paragraph{Language models as components of a decision system.} We use a language model for three tasks: reading a
part description, reading a document, and driving the system's tools. In none of them is its output trusted. The
instruction used for reading is optimised by reflective prompt evolution \citep{agrawal2025gepa}, which searches
over instructions with a language model as the mutation operator and keeps a Pareto front over examples; we
implement it locally, with a call budget and a held-out test. The concentration bound we use for the held-out
comparison is that of \citet{hoeffding1963}.

\paragraph{Forecasting intermittent arrivals.} Weekly counts of scrapped units are intermittent series, for which
the standard methods are those of \citet{croston1972}, \citet{syntetos2005} and \citet{teunter2011}. Zero-shot
time-series foundation models \citep{ansari2024chronos,das2024timesfm,woo2024,cohen2024toto,cohen2025time,mistral2026timestral}
are candidate competitors that need no training on the site's data. Forecast quality matters here through the
decision it supports, in the sense of \citet{elmachtoub2022}; Proposition~\ref{prop:fluid} bounds what a perfect
forecast could add.

\paragraph{Note on sources.} The descriptions above rest on the abstracts and bibliographic records of the cited
papers. Entries marked in the bibliography file as ``to verify'' were cited from memory. The claim that the
elements listed as contributions do not appear in this literature rests on these abstracts and on a limited number
of targeted searches; it is absence of evidence and must be confirmed on the full texts before submission.

\section{Setting and disposition model}\label{sec:model}

\subsection{Units, components, dispositions}

The model is written for abstract unit types, components, dispositions, access levels and resources, so that the
repair centre is one data instance and a production line another (Section~\ref{sec:discussion}). Each symbol has one
meaning throughout; periods are indexed by $k$, and the letter $t$ is kept for operator times.

Let $\cI$ be the set of types of scrapped unit and $\cJ_i$ the components of a unit of type $i$, with $n_{ij}\ge 1$
copies of component $j$ per unit. Each component receives a \emph{disposition} $d\in\cD$, where $0\in\cD$ means
``leave it in the unit''. A unit can be opened to an \emph{access level} $\ell\in\cL=\{0,1,\dots,L\}$, where $0$
means not opened; disposition $d$ requires level $\ell_d$, with $\ell_0=0$. A \emph{plan} for a unit of type $i$ is a
choice $s=(s_j)_{j\in\cJ_i}\in\cS_i=\cD^{\cJ_i}$ of one disposition per component; $s^0$ leaves every component.

The data of one period are: the number $N_i\ge 0$ of scrapped units of type $i$; the time $a_{i\ell}\ge 0$ to open a
unit of type $i$ to level $\ell$, with $a_{i0}=0$; the operator time $t_{jd}\ge 0$ per copy of $j$ under $d$, with
$t_{j0}=0$; the net value $\rho_{ijd}$ per copy before labour, with $\rho_{ij0}=0$; the cost $w_i\ge0$ of an operator
hour on type $i$; the hours $H\ge0$ available; and the fixed cost $K\ge 0$ of the line for the period. We write
$N^{c}_j=\sum_{i:\,j\in\cJ_i}N_i\,n_{ij}$ for the copies of component $j$ in the scrapped units of the period, all
types together.

A plan $s$ induces an access level, a labour content and a margin per unit:
\begin{equation}\label{eq:plan}
\ell(s)=\max_{j\in\cJ_i}\ell_{s_j},\qquad h_{is}=a_{i\,\ell(s)}+\sum_{j}n_{ij}\,t_{js_j},\qquad
g_{is}=\sum_{j}n_{ij}\,\rho_{ijs_j}-w_i\,h_{is}.
\end{equation}
The plan $s^0$ has $h_{is^0}=g_{is^0}=0$: the status quo is the reference and is worth zero.

\subsection{The instance of the repair centre}

There are three dispositions: reuse ($\R$, which needs a careful opening, level 2), sorted recycling ($\Sd$, for
which a rough opening suffices, level 1) and leave ($0$). Let $q_{ij}$ be the probability that an extracted copy is
reusable, $v_j$ the value of a reused copy, $\Lambda_j$ the value of the procurement lead time it avoids,
$\bar v_j=v_j+\Lambda_j$, $\kappa_j$ the transport cost, $m_j$ the mass of a copy, $p_j$ the price paid by the
recycler for its material class, and $\tau$ the cost per kilogram of the current bulk treatment (negative if the
current contractor pays). Then
\begin{equation}\label{eq:values}
\rho_{ij\R}=q_{ij}\,\bar v_j+(1-q_{ij})\,m_jp_j+\tau m_j-\kappa_j,\qquad \rho_{ij\Sd}=m_j\,(p_j+\tau),
\qquad t_{j\R}=t^{\mathrm{ext}}_j+t^{\mathrm{test}}_j,\qquad t_{j\Sd}=t^{\mathrm{sort}}_j .
\end{equation}
A copy that fails its test is sold by the kilogram, and every copy taken out of the unit avoids the bulk treatment.

\paragraph{Missing data are part of the model.} The class of a component fixes its price. When $p_j$ is unknown
and $m_j$ known, sorting $j$ is not available, a failed copy returns to the bulk stream, and
$\rho_{ij\R}=q_{ij}(\bar v_j+\tau m_j)-\kappa_j$. When $m_j$ is unknown, sorting is not available and
$\rho_{ij\R}=q_{ij}\bar v_j-\kappa_j$. We write $\out_j$ for the state ``the class of $j$ is not established''.
An unknown is never replaced by zero or by an average: the option that needs it disappears, and the terms left out
are reported. This convention is what makes the value of a missing fact well defined in Section~\ref{sec:knowledge}.

\subsection{The disposition problem}

Let $x_{is}\ge 0$ be the number of units of type $i$ processed with plan $s$. The disposition problem is
\begin{equation}\tag{P}\label{eq:P}
V(N,H)\;=\;\max_{x\ge 0}\ \sum_{i\in\cI}\sum_{s\in\cS_i}g_{is}\,x_{is}
\quad\text{s.t.}\quad \sum_{s\in\cS_i}x_{is}=N_i\ \ \forall i\in\cI,\qquad
\sum_{i\in\cI}\sum_{s\in\cS_i}h_{is}\,x_{is}\le H .
\end{equation}
The line is installed if $V(N,H)>K$, and the gain over the status quo is $V(N,H)-K$. Problem \eqref{eq:P} is always
feasible ($x_{is^0}=N_i$) and bounded. It is the linear relaxation of a multiple-choice knapsack problem: each type
is a class, each plan an item, the labour budget the knapsack. The relaxation is the natural model because $N_i$
counts many identical units; Remark~\ref{rem:integrality} bounds what integrality costs. The number of plans,
$|\cD|^{|\cJ_i|}$, is the difficulty: a unit of twenty components has more than three billion.

\begin{assumption}[Nested access]\label{as:nested}
$0=a_{i0}\le a_{i1}\le\dots\le a_{iL}$ for every $i$: opening a unit further never takes less time.
\end{assumption}
\begin{assumption}[Linearity]\label{as:linear}
Times and values are proportional to the number of units: no learning effect, no minimum lot.
\end{assumption}
\begin{assumption}[Independence across components]\label{as:indep}
The state of a component does not depend on the state of the others in the same unit; a type-level factor absorbs
the average effect of the unit having been scrapped.
\end{assumption}
Assumption~\ref{as:nested} is used only in Proposition~\ref{prop:pricing}, and is not needed for the two-level
instance (Remark~\ref{rem:twolevel}). Assumption~\ref{as:indep} is the weakest of the three; we return to it in
Section~\ref{sec:discussion}.

\section{Structure of the optimal disposition}\label{sec:structure}

All proofs are in Appendix~\ref{app:proofs}.

\subsection{Pricing a unit without enumerating its plans}

For $\mu\ge 0$ let $\varphi_i(\mu)=\max_{s\in\cS_i}(g_{is}-\mu h_{is})$ be the best \emph{reduced margin} of a unit
of type $i$ when an operator hour is priced at $\mu$ on top of its cost.

\begin{proposition}[Pricing by access level]\label{prop:pricing}
Under Assumption~\ref{as:nested}, for every $i$ and $\mu\ge0$, with $\omega_i=w_i+\mu$,
\begin{equation}\label{eq:pricing}
\varphi_i(\mu)\;=\;\max_{\ell\in\cL}\Big\{-\omega_i\,a_{i\ell}+\sum_{j\in\cJ_i}\ \max_{d:\,\ell_d\le\ell}\
n_{ij}\big(\rho_{ijd}-\omega_i\,t_{jd}\big)\Big\}.
\end{equation}
The pricing problem of a type is solved in $O(|\cJ_i|\,|\cD|\,L)$ operations.
\end{proposition}

The set-up couples the components of a unit, which is why plans cannot be priced component by component in general.
Proposition~\ref{prop:pricing} says that the coupling has only $L+1$ states: once the access level is fixed,
components are independent. Assumption~\ref{as:nested} is needed: when a rough opening takes longer than a careful
one, the right-hand side of \eqref{eq:pricing} can overstate $\varphi_i$.

\begin{remark}[Two levels, without Assumption~\ref{as:nested}]\label{rem:twolevel}
In the instance of Section~\ref{sec:model}, write $\Psi^{d}_j=n_{ij}(\rho_{ijd}-\omega_it_{jd})$ for
$d\in\{\R,\Sd\}$. Then
\[
\varphi_i(\mu)=\max\Big\{0,\ \ -\omega_ia_{i1}+\sum_j\pos{\Psi^{\Sd}_j},\ \
-\omega_ia_{i2}+\sum_j\max(\Psi^{\R}_j,\Psi^{\Sd}_j,0)-\min_j\pos{\max(\Psi^{\Sd}_j,0)-\Psi^{\R}_j}\Big\},
\]
where the second term is kept only if some $\Psi^{\Sd}_j\ge0$. The three terms are the best plans with no opening,
with sorting only, and with at least one reuse; in the last, each component takes its best disposition, and if none
chooses reuse it is imposed on the component for which it costs least. No ordering of $a_{i1}$ and $a_{i2}$ is
used. This is the routine of the implementation.
\end{remark}

\subsection{The value of an hour}

\begin{proposition}[Dual and value function]\label{prop:dual}
\begin{equation}\label{eq:dual}
V(N,H)\;=\;\min_{\mu\ge0}\Big\{\mu H+\sum_{i\in\cI}N_i\,\varphi_i(\mu)\Big\}.
\end{equation}
$V$ is concave, piecewise linear and positively homogeneous of degree one in $(N,H)$, and nondecreasing in each
argument. If $\mu^\ast$ attains the minimum, then $V(N,H')\le V(N,H)+\mu^\ast(H'-H)$ for every $H'\ge0$. Any optimal
$\mu^\ast$ lies between the value of one more operator hour and the loss from one hour less, which coincide except
at the finitely many breakpoints of $V$; and $\mu^\ast=0$ is optimal when the labour budget is not exhausted.
\end{proposition}

The dual objective is convex and piecewise linear in $\mu$, with subgradient $H-\sum_iN_ih_{is_i(\mu)}$ at any
selection $s_i(\mu)$ of best plans. The implementation locates $\mu^\ast$ by bisection on the sign of this
subgradient, each evaluation costing one pricing per type by Proposition~\ref{prop:pricing}, and recovers a primal
solution by Proposition~\ref{prop:split}.

\begin{proposition}[Few split types, efficient plans]\label{prop:split}
Problem \eqref{eq:P} has an optimal solution in which every type with $N_i>0$ uses one plan, except at most one
type, which uses two. In any optimal solution, every plan in use attains $\varphi_i(\mu^\ast)$ for an optimal
multiplier $\mu^\ast$; in particular it lies on the upper concave envelope of the points
$\{(h_{is},g_{is})\}_{s\in\cS_i}$, and the two plans of a split type lie on a common edge of that envelope, of
slope $\mu^\ast$. With $M$ coupling constraints in place of the single labour budget, a feasible programme has an
optimal solution in which at most $M$ types use more than one plan.
\end{proposition}

For a single budget this is the classical structure of the linear multiple-choice knapsack
\citep{sinha1979,zemel1980}; we state it for $M$ constraints because circularity targets add constraints.

\begin{remark}[Integrality]\label{rem:integrality}
With a single budget, rounding the split type down to whole units and sending the remainder to the cheaper of its
two plans keeps the solution feasible and loses less than the margin difference of one unit. The argument holds for
$\le$ constraints only; a lower bound on an outcome can be violated by rounding down.
\end{remark}

\paragraph{Managerial reading.} At the optimum a single number, the value of an hour $\mu^\ast$, coordinates all
decisions: every type is processed with the plan that is best when labour is charged $w_i+\mu^\ast$, and at most
one type is processed in two ways, because the hours run out part-way through it. The same number tells the manager
what an additional operator hour is worth, which is the first question asked once a line exists.

\subsection{A threshold rule for the bench}

Consider a unit that is opened carefully anyway. The choice is then separable across components.

\begin{proposition}[Yield threshold]\label{prop:threshold}
In the instance of Section~\ref{sec:model}, let $\omega=w_i+\mu$ and assume $\bar v_j>m_jp_j$. In a unit opened to
level 2, reuse is at least as good as both sorting and leaving component $j$ if and only if
\begin{equation}\label{eq:threshold}
q_{ij}\;\ge\;q^{\dagger}_{ij}(\mu)\;=\;\frac{\kappa_j+\omega\,t_{j\R}-\min\{\omega\,t_{j\Sd},\ m_j(p_j+\tau)\}}
{\bar v_j-m_jp_j}.
\end{equation}
When $t_{j\R}\ge t_{j\Sd}$ the threshold is nondecreasing in $\mu$.
\end{proposition}

The rule is what an operator can apply: \emph{extract and test component $j$ when its observed good rate is above
$q^\dagger_{ij}$}. It also says how scarcity of labour changes behaviour: the scarcer the hours, the higher the
yield a component must show. Both hypotheses are needed: the comparison reverses when a good copy is worth less
reused than as material, and the threshold can decrease in $\mu$ when testing is faster than sorting. The threshold
is conditional on the unit being opened carefully; whether it is opened at all is decided by
Proposition~\ref{prop:pricing}, where the set-up is weighed against the sum of the components' reduced values.

\subsection{Circularity objectives}

Three outcomes are linear in $x$ and are reported for every solution: the mass diverted from the bulk stream,
$\sum_{i,s}b_{is}x_{is}$ with $b_{is}$ the mass taken out of a unit under plan $s$; the mass of critical raw
materials sent to a recovery route; and the number of parts returned to service, $\sum_{i,s}e_{is}x_{is}$ with
$e_{is}=\sum_{j:s_j=\R}n_{ij}q_{ij}$.

\begin{corollary}[Frontier]\label{cor:frontier}
Adding a weighted linear outcome to the objective keeps at most one split type. Imposing $M$ lower bounds on linear
outcomes gives $M+1$ coupling constraints; when the programme is feasible it has an optimal solution with at most
$M+1$ split types.
\end{corollary}

The frontier between margin and parts returned to service answers a question that a single-objective optimum
cannot: how much margin does each additional part returned to service cost? Returning parts to service lowers the
cost of repair and of refurbished systems; whether this improves access to imaging for hospitals with small budgets
is an empirical question outside the model, and we make no claim about it.

\section{Unknown yields and irregular arrivals}\label{sec:learning}

\subsection{A prior from repair records, corrected by teardown}

The yield $q_{ij}$ is unknown before any unit has been torn down. Repair records give a starting point. Let $\nu_i$
be the number of repairs of type $i$ on record, $\mathrm{cons}_{ij}$ the copies of $j$ consumed in them, and
$\eta_j$ the share of replacements made as a precaution. The probability that a copy is defective in a
\emph{repaired} unit is estimated by $\hat\delta_{ij}=(1-\eta_j)\,\mathrm{cons}_{ij}/(\nu_in_{ij})$, and the prior
mean of the yield is
\begin{equation}\label{eq:prior}
q^0_{ij}=\big[1-\gamma_i\,\hat\delta_{ij}\big]_{[0,1]}\,(1-\chi_j)\,\vartheta_j,
\end{equation}
with $\gamma_i\ge1$ a factor for scrapped units being in worse condition than repaired ones, $\chi_j$ the breakage
at extraction and $\vartheta_j$ the share of sound copies that pass the test. Teardown audits and routine tests then
give binomial observations: of $z$ copies tested, $y$ are good. With prior
$\mathrm{Beta}(q^0\varpi,(1-q^0)\varpi)$, where $\varpi$ is the weight of the prior in equivalent tested copies, the
posterior is $\mathrm{Beta}(q^0\varpi+y,\ (1-q^0)\varpi+z-y)$.

The prior is a modelling assumption with three weaknesses, each of which biases it in a known direction: preventive
replacement inflates $\hat\delta$; $\gamma_i$ is not identified without teardown data; and repairs that consumed
nothing are missing from $\nu_i$. The prior is there to start; the audit is there to correct it.

\subsection{Why a greedy rule fails, and posterior sampling}

The state of a copy is observed only if it is extracted and tested. A rule that follows the posterior mean never
tests a component whose prior mean is below the threshold of Proposition~\ref{prop:threshold}, and therefore never
corrects a pessimistic prior: the error is permanent because the data that would reveal it are never produced. The
problem is a combinatorial semi-bandit: each period the decision is a feasible point of \eqref{eq:P}, whose labour
budget is renewed, the reward is affine in the unknown yields, and the outcome of each tested copy is observed.

\begin{center}
\fbox{\parbox{0.92\textwidth}{\textbf{Algorithm TS-LP.} In each period $k$:
(1) draw $\tilde q_{ij}$ from the current posterior of every pair $(i,j)$;
(2) solve \eqref{eq:P} with $\tilde q$ and the arrivals of the period, by Propositions~\ref{prop:pricing}
to~\ref{prop:split};
(3) process the units accordingly;
(4) update the posteriors with the test outcomes.}}
\end{center}

With $x_k$ the plan of period $k$ and $N_k$ its arrivals, regret is
$\sum_{k\le T}\big(V(N_k,H;q)-\E[g(q)^\top x_k]\big)$. We evaluate it numerically
(Section~\ref{sec:experiments}) and claim no bound, for the reason given in Section~\ref{sec:literature}.

\paragraph{Scoring the estimate.} The posterior is scored prequentially: each audit batch is predicted from what
was known before it and only then used for updating. We report the Brier score and the logarithmic loss for grouped
binomial outcomes, the absolute error on the batch rate, and the coverage of the 90\% predictive interval. Two
references must be beaten: the prior alone, and the raw observed frequency.

\paragraph{Reused parts must also last.} A reusable part must survive once installed. The field file available
records only early failures within a window of $W$ days after installation, in whole days, with no survivors. From
such data the Weibull shape is estimable by a right-truncated discrete likelihood, in which the probability of a
failure on day $\iota$ is $[F(\iota+1)-F(\iota)]/F(W+1)$, while the scale is not identifiable; early-failure rates
of repaired or harvested parts relative to new ones are estimable when the number installed is known. The expected
cost of an early failure enters the model through the value $v_j$.

\subsection{What buffering, forecasting or reinforcement learning can gain}

Let $A_k=(A_{ik})_i$ be the arrivals of period $k$, with finite means and
$\bar A=\lim_T\frac1T\sum_{k\le T}\E A_k$, and let $H$ hours be available in every period. A \emph{policy} chooses in
each period, from the history, how many units of each type to process with each plan among the units on hand,
within $H$ hours; units not processed can be kept, within any storage limit, or released to the status quo. The
\emph{myopic} policy keeps no stock and solves \eqref{eq:P} on each period's arrivals.

\begin{proposition}[Fluid bound and variability gap]\label{prop:fluid}
Assume the initial stock is empty, the margins $g$ are known and $\bar A$ exists. For every policy, with $x_k$ the
plan it applies in period $k$,
\begin{equation}\label{eq:fluid}
\limsup_{T\to\infty}\frac1T\,\E\Big[\sum_{k=1}^{T}g^\top x_k\Big]\;\le\;V(\bar A,H),
\end{equation}
and the myopic policy satisfies $\frac1T\sum_{k\le T}\E\,V(A_k,H)\le V\big(\frac1T\sum_{k\le T}\E A_k,\,H\big)$.
Hence
\begin{equation}\label{eq:gamma}
\Gamma\;=\;V(\bar A,H)-\liminf_{T}\frac1T\sum_{k\le T}\E\,V(A_k,H)\;\ge\;0
\end{equation}
is an upper bound on what any policy that stores units, forecasts arrivals or is trained by reinforcement learning
can gain per period over the myopic policy. Pathwise, for every realised sequence and every policy, including one
that knows the future, $\sum_{k\le T}g^\top x_k\le T\,V\big(\tfrac1T\sum_{k\le T}A_k,\,H\big)$.
\end{proposition}

The proposition turns a design question into a measurement. The sample version of $\Gamma$, the model run on the
average flow minus the model replayed period by period, is computable from the history. It is an exact hindsight
bound for the observed sequence, and it estimates $\Gamma$ with a downward bias on short histories, since
$\E\,V(\text{sample mean},H)\le V(\text{mean},H)$. If it is small relative to $V$, no buffer, forecast or learned
policy can add much, and the myopic rule is enough. If it is large, the candidates are, in increasing order of
complexity, a buffer with a base-stock rule, a rolling-horizon version of \eqref{eq:P} fed with forecasts, a
stochastic programme on forecast scenarios, and a policy trained by reinforcement learning on a simulator; each is
judged by the share of $\Gamma$ it recovers. The bound is not tight in general, since a finite buffer cannot smooth
arbitrarily long droughts, and it takes the margins as known: it does not bound the value of learning the yields.

\subsection{Forecasting arrivals}\label{sec:forecast}

Weekly counts of scrapped units per type are short, low-count, intermittent series. They have, however, a known
driver. A unit is scrapped when a unit received for repair is found beyond economic repair. With $E_{ik}$ the units
of type $i$ that reached a decision in period $k$ and $\sigma_i$ the scrap rate,
\begin{equation}\label{eq:structural}
A_{ik}\mid E_{ik},\sigma_i\sim\mathrm{Binomial}(E_{ik},\sigma_i),\qquad\sigma_i\sim\mathrm{Beta}(\alpha_i,\beta_i).
\end{equation}
The prior $(\alpha_i,\beta_i)$ of a type is estimated from the \emph{other} types, by weighted moments from which
the expected binomial sampling variance is removed, so that sparse types borrow strength and the data of a type are
not used twice. The inflow $E_{ik}$ is a denser series, driven by the installed base, and is forecast separately;
the predictive distribution of next period's arrivals is Beta--binomial given a forecast of its inflow.

The structural model is compared with the historical mean, the methods of \citet{croston1972},
\citet{syntetos2005} and \citet{teunter2011}, and a pooled Gamma--Poisson model. Zero-shot foundation models can be
plugged in through an adapter on a machine that holds their weights; none was run for this paper. The protocol is a
rolling origin with expanding window, all methods being scored on the same (type, period) pairs. Point accuracy is
the root mean squared error relative to that of the historical mean; the absolute error is not used, because it is
minimised by the median, which is zero as soon as more than half of the periods are empty. Probabilistic accuracy
is the pinball loss on a grid of quantiles and the coverage of the 80\% interval.

\section{Knowledge acquisition: what is worth knowing}\label{sec:knowledge}

\subsection{The missing fact and the means to obtain it}

The recycling price $p_j$, the regulatory route and the content in critical raw materials all depend on the
\emph{material class} of component $j$, which the bill of materials does not give. The class can be obtained by
means of very different cost and reliability: a rule on the description; analogy with a part already validated; an
earlier classification; a reading of the internal description by a language model; a manufacturer's material
declaration or a data sheet retrieved from the web; a reading of such a document by a language model; a check of
that reading; a human expert. With thousands of components, most of them light or rare, the question is not how to
classify a component but which components deserve which effort.

Let $\cC$ be the material classes of the recycler's reference list and $\bC=\cC\cup\{\out\}$, where $\out$ stands
for ``a class outside the list''; let $n=|\bC|$. Let $C_j\in\bC$ be the unknown class of component $j$, $p(c)$ the
price of class $c\in\cC$, and
\[
v(c)=\pos{p(c)+\tau}\ \ (c\in\cC),\qquad v(\out)=0,
\]
the most a kilogram of class $c$ can earn over the bulk stream. Quantities in this section are per year, the period
over which an established class is taken to stay valid; $M_j$ is the annual mass of the copies of $j$ in the streams
the decision covers.

\subsection{A belief over the class, with sources of learned reliability}\label{sec:belief}

The belief $\pi_j$ over $\bC$ starts uniform. A piece of evidence is a triple (source type $u$, class $c$, sign) of
one of two kinds. An \emph{assertion} names a class among all: a source of reliability $r$ names the true class with
probability $r$ and otherwise one of the $n-1$ others, uniformly. A \emph{verdict} answers yes or no to the question
``is it $c$?'': a source of reliability $r$ answers correctly with probability $r$, whatever the true class. Pieces
of the same source type share one unknown reliability $r_u\sim\mathrm{Beta}(a_u,b_u)$; given the truth and the
reliabilities, pieces of independent origin are independent. For a candidate class $c$, let $s_u(c)$ and $f_u(c)$
be the numbers of pieces of type $u$ that would be right and wrong if $C_j=c$, and $e_u(c)$ the number of wrong
assertions among them. Integrating over the reliabilities,
\begin{equation}\label{eq:fusion}
\Pr(C_j=c\mid\text{evidence})\;\propto\;\prod_{u}\frac{\Bfun\big(a_u+s_u(c),\,b_u+f_u(c)\big)}{\Bfun(a_u,b_u)}\,
(n-1)^{-e_u(c)},
\end{equation}
where $\Bfun$ is the Beta function. With one piece per source type the factor is the mean reliability when the piece
is right and its complement when it is wrong; with several pieces of one type it is less confident than the product
of means, because pieces that share an unknown reliability are not independent. For assertions this is the
observation model of \citet{dawid1979} with one accuracy per source type and uniform errors.

Four design decisions follow from treating every automatic source as fallible.
\begin{itemize}
\item \emph{A language model is a source, never ground truth.} Its readings and verdicts enter \eqref{eq:fusion} with
  their own reliability. Ground truth is a human validation.
\item \emph{Trust is earned.} With the declared priors, a document read by a model and checked by the same model
  does not reach the acceptance threshold; it does once validations have shown that this kind of reading is right
  often enough.
\item \emph{Independence is enforced by provenance.} Ten web pages reproducing one data sheet are one piece of
  evidence. A model that re-reads its own reading adds no source: a ``yes'' turns the reading into a confirmed
  reading, a ``no'' withdraws it. Two readers of the same text that agree count once, as the more reliable of the
  two.
\item \emph{Unknown stays unknown.} A component whose class is not established has no price: in \eqref{eq:P} it is
  in the state $\out_j$, where sorting is unavailable. It is never valued at an average price.
\end{itemize}

\subsection{The value of a fact}

Two cases must be kept apart. For a parameter on which the decision maker can act at its mean (a yield, a value,
the price of an established class known with error), the reference is the decision taken at the mean. For the class
of a component there is no decision at the mean class: without a class the component cannot be sorted.

\begin{proposition}[Value of information on a parameter]\label{prop:evpi}
For fixed $(N,H)$, $V$ is a convex function of the vector of values $(\rho_{ijd})$, hence of every parameter
$\theta$ that enters these values affinely. If $\theta$ has distribution $F$ and mean $\bar\theta$, the expected
value of perfect information $\EVPI=\E_F[V(\theta)]-V(\bar\theta)$ is nonnegative, and it is zero if and only if a
solution optimal at $\bar\theta$ is optimal for $F$-almost every $\theta$. If the slopes of $V$ in $\theta$ lie in
an interval of length $\Delta$, then
\begin{equation}\label{eq:evpi}
\EVPI\;\le\;\tfrac12\,\Delta\;\E_F|\theta-\bar\theta|,
\end{equation}
and the constant $\tfrac12$ cannot be improved. In particular
$\EVPI\le\tfrac12\,N^{c}_jm_j\,\E_F|p_j-\bar p_j|$ for the price of component $j$, and
$\EVPI\le\tfrac12\,N_in_{ij}\,|\lambda_{ij}|\,\E_F|q_{ij}-\bar q_{ij}|$ for a yield, where $\lambda_{ij}$ is the
coefficient of $q_{ij}$ in $\rho_{ij\R}$.
\end{proposition}

\begin{proposition}[Value of establishing a class]\label{prop:class}
In the instance of Section~\ref{sec:model}, let the class of component $j$ be unknown with belief $\pi$. Let
$V(\out_j)$ be the value of \eqref{eq:P} when $j$ has no price and $V(p)$ its value when the price of $j$ is $p$.
For every price $p$, $V(p)-V(\out_j)\le N^{c}_jm_j\pos{p+\tau}$, and therefore
\begin{equation}\label{eq:class}
\E_\pi\big[V(p(C_j))\big]-V(\out_j)\;\le\;N^{c}_j\,m_j\ \E_\pi\big[v(C_j)\big].
\end{equation}
The left-hand side is nonnegative when $p(c)+\tau\ge0$ for every class of positive probability.
\end{proposition}

The left-hand side of \eqref{eq:class} can be negative when a possible class costs more to recycle than the bulk
stream: establishing the class then reveals an obligation, not an opportunity, and the decision to establish it is
not an economic one. The bound extends unchanged to copies of $j$ that leave the site by another stream (defective
copies replaced in repairs), since sorting a copy earns at most $m_jv(C_j)$; this gives the annual mass $M_j$.

\subsection{The acquisition chain and its budget}

\begin{definition}[Stake]\label{def:stake}
The \emph{stake} of component $j$ under belief $\pi$ is $e_j(\pi)=M_j\sum_{c\in\bC}\pi(c)\,v(c)$, and its
\emph{ceiling} is $\bar e_j=M_j\max_{c\in\cC}v(c)$. The stake is unknown when $M_j$ is unknown or when a class of
non-negligible probability has no price.
\end{definition}

By Proposition~\ref{prop:class}, the stake bounds from above what establishing the class of $j$ can earn in a year.
The chain uses it as a budget. Let $S_j$ be the cumulative cost of the paid acquisitions already made for $j$, over
all runs; it is kept in a journal. Each paid action $a$ has a cost $\phi_a>0$.

\begin{center}
\fbox{\parbox{0.92\textwidth}{\textbf{Algorithm 1 (acquisition chain for component $j$).}
\begin{enumerate}\setlength{\itemsep}{0pt}
\item Record the free evidence: rule on the description, analogy with a validated part, earlier classification.
  Update $\pi$. If a human validation exists, stop.
\item For each paid action $a$, in the order: reading of the internal description by the private model; document
  retrieval; reading of each retrieved document by a model; verdict of a model on the retained class (attempted
  only if a ``yes'' would bring the class to acceptance):
  \emph{if the class is not accepted, the action applies, and $e_j(\pi)\ge S_j+\phi_a$, perform $a$, set
  $S_j\leftarrow S_j+\phi_a$ and update $\pi$.} An action already paid for is not repeated; an action that fails
  for a technical reason is recorded, not charged, and retried at the next run.
\item \emph{Accept} the most probable class $\hat c\ne\out$ without review if $\pi(\hat c)\ge\theta$ and the value at
  risk $R_j=(1-\pi(\hat c))\,M_j\max_{c}v(c)$ is below the cost $\phi_{\mathrm{rev}}$ of a review. Otherwise send
  the component to the human queue, ranked by stake, if $e_j(\pi)\ge\phi_{\mathrm{rev}}$; else leave it.
\end{enumerate}
Components that are not sent to a person (accepted, or left) are audited by random sampling; queue decisions and audit
decisions return to the knowledge base as
validations, tagged with their reason.}}
\end{center}

\begin{lemma}[Budget invariant]\label{lem:budget}
For every component $j$ and at every time, $S_j\le\bar e_j$. Consequently the total spending on a bill of materials
is at most $\sum_j\bar e_j$; a component with $\bar e_j<\min_a\phi_a$ triggers no paid action, in particular no call
to a language model; and neither does a component whose stake is unknown.
\end{lemma}

The lemma is elementary, and it is the property that makes the layer safe to run unattended on a bill of materials
of any size: spending is bounded by value, component by component, whatever the sources answer.

\begin{assumption}[No gain from error]\label{as:noerror}
Whatever the class under which component $j$ is handled, if its true class is $c$ the realised value of the
programme is at most $\max\{V(p(c)),V(\out_j)\}$.
\end{assumption}

\begin{proposition}[The screen is sound]\label{prop:screen}
Under Assumption~\ref{as:noerror}, consider any acquisition procedure for component $j$ started at belief $\pi$: an
adaptive sequence of actions with total cost $\Phi$, followed by a decision on the class. Let $G$ be its gain over
leaving the class unestablished. Then $\E_\pi[G]\le e_j(\pi)$. Consequently:
\begin{enumerate}\setlength{\itemsep}{0pt}
\item[(i)] a procedure with $\E_\pi[\Phi]>e_j(\pi)$ has negative expected net value;
\item[(ii)] on a component for which nothing has been spent, the test $e_j(\pi)\ge\phi_a$ excludes no action whose
  expected net value is positive;
\item[(iii)] under the evidence model of Section~\ref{sec:belief}, the stake $e_j(\pi_t)$ along the chain is a
  martingale: no action is justified by the expectation that it will raise the stake.
\end{enumerate}
\end{proposition}

Once something has been spent, the test of Algorithm~1 charges past spending against the stake. It is then stricter
than (ii): it is a budget rule, which can forgo an action whose marginal value exceeds its cost, in exchange for
Lemma~\ref{lem:budget}. We accept this conservatism because the stake is itself an upper bound, not the value. The
chain follows a fixed cheapest-first order; ranking actions by expected value per unit cost is not implemented.

\begin{assumption}[A wrong class loses at most the right one]\label{as:loss}
Handling component $j$ under a class other than its true class $c$ loses, relative to handling it under $c$, at most
$M_jv(c)$ per year; and a human review reveals the true class at cost $\phi_{\mathrm{rev}}$.
\end{assumption}

\begin{lemma}[Acceptance]\label{lem:accept}
Under Assumption~\ref{as:loss}, accepting $\hat c$ under belief $\pi$ has an expected loss of at most
$M_j\sum_{c\ne\hat c}\pi(c)v(c)\le R_j$. Hence when $R_j<\phi_{\mathrm{rev}}$, accepting costs less in expectation
than reviewing. If moreover beliefs are calibrated, the expected share of wrong classes among accepted components
is at most $1-\theta$.
\end{lemma}

The two conditions of the acceptance rule do different work. The value at risk is the economic condition. The
probability threshold is not: without it, light components would be accepted on thin evidence, and since accepted
components are those the audit samples to learn the reliability of sources, the threshold protects the measurement.
Assumption~\ref{as:loss} ignores the penalty a recycler may charge for a contaminated stream, which would enter as a
larger value at risk.

\subsection{Why the fusion converges, and how fast}

\begin{proposition}[Identification from a source of known reliability]\label{prop:fusion}
Let $K$ independent assertions come from a source of known reliability $r$, with $n\ge2$ possible classes and a
uniform prior. Let $q=(1-r)/(n-1)$ and $k_c$ the number of assertions naming $c$. The posterior is proportional to
$(r/q)^{k_c}$: the most probable class is the most named when $r>1/n$ and the least named when $r<1/n$; when
$r=1/n$ the posterior equals the prior whatever $K$. For $r\ne1/n$, the most probable class $\hat c_K$ satisfies
\begin{equation}\label{eq:rate}
\Pr\big(\hat c_K\ne C_j\big)\;\le\;(n-1)\,\varrho^{K},\qquad\varrho=1-\big(\sqrt r-\sqrt q\,\big)^2<1,
\end{equation}
where a tie counts as an error. Hence $K\ge\ln((n-1)/\delta)/\ln(1/\varrho)$ assertions suffice for an error
probability of at most $\delta$, and the posterior probability of the true class tends to one almost surely.
\end{proposition}

Two readings matter in practice. First, what makes a source useful is not that it is right but that its reliability
is \emph{known and different from chance}: a source that is wrong more often than chance identifies the class as
surely as one that is right. This is why the layer invests in measuring reliabilities. Second, the rate depends on
$(\sqrt r-\sqrt q)^2$, so a few sources of high reliability are worth many mediocre ones: with \SynFusionClasses{} classes, a source
of reliability $0.9$ needs \SynFusionKhaut{} independent assertions for an error bound of 2\%, and a source of
reliability $0.6$ needs \SynFusionKbas.
The implemented rule \eqref{eq:fusion} integrates over an unknown reliability and reduces to the rule of the
proposition as validations accumulate.

\subsection{Measuring reliabilities from validations that are not a random sample}

Each $r_u$ starts from a weak declared prior, equivalent to four observations. A human validation of a component
shows which of its pieces of evidence were right. But validated components are not a random sample: those sent to
review are the ones on which the sources disagree, where at least one is wrong.

\begin{proposition}[Stratified estimation]\label{prop:strata}
Let the pieces of evidence of a source type be partitioned into a stratum $\mathsf A$, on components the chain does
not send to a person (accepted without review, or left because their class would earn less than a review costs),
and a stratum $\mathsf Q$, on components it sends to review, of sizes $n_{\mathsf A},n_{\mathsf Q}$
and accuracies $r_{\mathsf A},r_{\mathsf Q}$; let $n=n_{\mathsf A}+n_{\mathsf Q}$ and
$r=(n_{\mathsf A}r_{\mathsf A}+n_{\mathsf Q}r_{\mathsf Q})/n$. Suppose $k_{\mathsf A}$ pieces of $\mathsf A$ and
$k_{\mathsf Q}$ pieces of $\mathsf Q$ are validated, drawn by simple random sampling without replacement in each
stratum, of which $y_{\mathsf A}$ and $y_{\mathsf Q}$ are right; let $k=k_{\mathsf A}+k_{\mathsf Q}$.
\begin{enumerate}\setlength{\itemsep}{0pt}
\item[(i)] The stratified estimator $\hat r=(n_{\mathsf A}\,y_{\mathsf A}/k_{\mathsf A}+n_{\mathsf Q}\,y_{\mathsf Q}/k_{\mathsf Q})/n$
  is unbiased for $r$.
\item[(ii)] The pooled estimator $(y_{\mathsf A}+y_{\mathsf Q})/k$ has bias
  $(r_{\mathsf A}-r_{\mathsf Q})\,(k_{\mathsf A}/k-n_{\mathsf A}/n)$, which is negative when
  $r_{\mathsf Q}<r_{\mathsf A}$ and the review stratum is over-represented among validations.
\item[(iii)] With weights $n_s-k_s$ in place of $n_s$, the estimator is unbiased for the expected accuracy of the
  pieces \emph{not yet validated}, which is the quantity the fusion of the remaining components needs.
\end{enumerate}
\end{proposition}

The implementation uses (iii), with each stratum mean shrunk towards the declared prior with the weight of four
observations. Validations made without regard to what the sources say update the prior of both strata. Validations
of components picked by the reviewer set the class but do not judge the sources: they are neither random nor
exhaustive. The definition of the strata matters. The audit must draw among \emph{everything} the chain shows to
nobody, including the components it leaves for want of stake: an earlier version of the system put those in the
review stratum, where only contested, high-stake components are ever validated, and so lent the errors found there
to all the components nobody looks at. A closed-loop simulation exposed the defect (it is corrected in the code and
in this statement). The estimator is unbiased only as far as the audit is random and the components sent to review
are all reviewed; a queue reviewed from the top is not a random sample of the queue. Unbiasedness also has a price
in variance: in our simulations (not reported here) the stratified estimate removed the downward bias of the pooled
one without improving its mean absolute error.

\subsection{Confidentiality}\label{sec:confidentiality}

A company part number and the description of an in-house design are not public. An outward query contains only a
manufacturer's name, a manufacturer's part number and a generic component type, and it is issued only if the
component is not flagged as designed in-house (any unrecognised value of the flag counts as in-house), a
manufacturer is named and is not the company itself, and the part number has the form of a manufacturer's reference
and is not the company part number of any component of the bill of materials. Queries are built from
bill-of-materials fields \emph{before} any model call, so that nothing a model answers, including under the
influence of a hostile document, can enter a request. The language model is self-hosted: only a model on the local
machine, or on a machine of the owner's network that the owner has declared, may read internal descriptions, which
makes it the one automatic source for in-house designs. Any other server is refused unless explicitly accepted, and
then receives public documents only.

\section{The system: agents on a self-hosted open model}\label{sec:system}

\subsection{Architecture}

Figure~\ref{fig:architecture} shows the system. Operational records and a parameter workbook feed the decision
models of Sections~\ref{sec:model} to~\ref{sec:learning}. The workbook holds every parameter as a range with its
provenance; it is never rewritten by the software, and an unknown value stays unknown through every computation.
The value-of-information layer turns the disposition model into a stake for each component; the acquisition chain
spends within that stake and writes to a knowledge base; a review page presents the queue and an audit sample to a
human, whose decisions return as validations. All of this runs as a pipeline of seven steps whose outputs are
stored with a manifest. On top of it, three agents drive the pipeline and write the decision note.

\begin{figure}[t]
\centering
\includegraphics[width=\textwidth]{\genere/fig_architecture.pdf}
\caption{Architecture. The disposition model prices labour and, through the value of information, gives each
component a stake; the acquisition chain spends within it, on the components of the bill of materials; human
validations feed the reliability of sources and the search for a better reading instruction, which the chain then
uses. Agents call these modules through a closed list of tools.}\label{fig:architecture}
\end{figure}

\paragraph{The knowledge base.} Claims are stored with their source type, origin, sign, kind (assertion or
verdict) and the text they rest on; validations with their author and reason; and every acquisition with its cost
and outcome. The store is domain-neutral: it takes any subject and attribute, and only the acquisition chain is
specific to material composition. Because the journal records what was paid for, a search or a reading is never paid
twice, and a run interrupted by a failure resumes where it stopped.

\subsection{A private model, by construction}

The language model is an open-weight model served on hardware the data owner controls, through the
chat-completions interface that the common local servers share. The software classifies the server before using
it: \emph{this machine}; \emph{a declared trusted server} on the owner's network, named in a settings file;
\emph{any other server}, which is refused unless explicitly accepted and then never receives internal text. Only a
model of the first two kinds is \emph{private}. Reading internal descriptions, running the agents and optimising
instructions all require a private model and raise an error otherwise. The model is asked for structured output
against a schema; answers are parsed defensively, an answer outside the list of classes is discarded, and a failure
of the server is recorded and never charged. An offline diagnostic reports the packages, the graphics memory, the
local servers that answer and a model size that fits, and checks the integrity of the installed files against a
manifest of hashes; it is the first thing the launcher runs, so that a locked-down workstation needs no
back-and-forth to set up.

\subsection{Three agents and a closed list of tools}

The agents follow the pattern of a tool-using assistant, with the hosted model replaced by the private one. The
model computes nothing itself: it chooses tools, reads their results, decides what to do next, then writes.
\begin{itemize}
\item The \emph{pilot} receives the list of tools and answers, at each turn, with one JSON object naming a tool and
  its arguments. The tools are functions of the system: run the diagnostic; run the pipeline up to a step; read a
  result table; read the decision indicators; read the limits and the missing data; run the acquisition chain; read
  the review queue. The number of turns is bounded and each turn is logged.
\item The \emph{writer} receives the pilot's findings and the evidence (the text returned by the tools) and writes
  the decision note for the manager.
\item The \emph{reviewer} reads the note against the evidence and records its reservations, which are appended to
  the note.
\end{itemize}
Two properties do not depend on the model behaving well.

\begin{lemma}[Containment and grounding]\label{lem:grounding}
Whatever the model outputs: (i) the effects of an agent run are those of the listed tools, since any other request
is answered by an error message and executes nothing; (ii) every numeral of the issued note that has decimals or is
at least ten either matches a numeral of the evidence (exactly, at the note's rounding, or as the percentage of a
proportion) or is listed in a warning placed at the head of the note; (iii) if the model fails at any point, the
run terminates with a fixed plan and a template note computed from the same results.
\end{lemma}

The check in (ii) is a program, run after the last model call; a note whose first version fails it is rewritten
once with the list of offending numbers, and what still fails is flagged, never silently passed. The lemma
guarantees that a number in the note \emph{exists} in the results. It does not guarantee that the sentence around
it describes the right quantity; that is what the reviewer and, finally, the reader are for. Without a model the
fixed plan runs the same tools in a set order: the proof of concept runs on any workstation, and gains in
readability, not in correctness, when a model is connected.

\subsection{Improving the reading instruction from validations}

The instruction given to the private model for reading descriptions is learned from the human validations by
reflective prompt evolution \citep{agrawal2025gepa}. Candidate instructions are scored per example; a candidate is
drawn among those that are best on at least one example of a selection set; a reflection model reads its errors on
a small batch and proposes a rewrite; the rewrite is kept only if it does better than its parent on that batch.
An abstention is scored one half, so that the model gains by answering only when it is more likely right than
wrong. Our implementation is local and independent of the reference implementation, and omits its merge step. It
adds three guards. A rewrite that contains a part number or copies a description is rejected, so that the
instruction cannot memorise the validated parts. The search has a hard budget $B$ of reading calls and caches every
(instruction, example) pair. And the validated parts are split at random into a batch set, a selection set and a
test set that the search never sees.

\begin{lemma}[Held-out adoption]\label{lem:heldout}
Suppose the validated parts are independent draws from a population of parts, and let the test set contain $n_T$ of
them that the search never sees, with scores in $[0,1]$. Let $\hat\Delta$ be the difference between the mean test
scores of the returned instruction and of the starting one, and $\Delta$ the difference of their expected scores on
a part of that population. Then $\Pr(\hat\Delta-\Delta\ge\epsilon)\le\exp(-n_T\epsilon^2/2)$ for every
$\epsilon>0$. In particular the rule ``adopt if $\hat\Delta\ge0$'' adopts an instruction that is worse than the
starting one by more than $\epsilon$ with probability at most $\exp(-n_T\epsilon^2/2)$. The search makes at most
$B$ reading calls and the test at most $2n_T$.
\end{lemma}

The guarantee holds for one search on one test set, and with a few dozen validated parts it is weak; the system
therefore also reports an exact paired test and states when the measured difference is compatible with chance. The
population is that of validated parts, which over-represents contested and high-stake components: the gain on the
whole bill of materials can differ from the gain measured.
No weight is changed and no record leaves the owner's machine.

\subsection{From one workstation to a virtual machine}\label{sec:deployment}

The same package runs at three scales, and what changes between them is one settings file.
\begin{enumerate}\setlength{\itemsep}{0pt}
\item \emph{Proof of concept.} Any workstation with a Python interpreter, without a graphics card and without a
  model: the fixed plan computes every result and writes the template note. With a small open model served
  locally, the agents take over.
\item \emph{Scale-up.} A workstation with a graphics card serves a larger model, on the same machine or declared
  as a trusted server for other workstations of the site. Reading the descriptions of a whole bill of materials
  runs several requests in parallel; the budget of Lemma~\ref{lem:budget} bounds their number.
\item \emph{Virtual machine.} The package, its settings and a model server run in a container on a virtual machine
  of the owner's infrastructure; the knowledge base is one file, and the manifest of each run makes results
  comparable across machines.
\end{enumerate}
\apres{the hardware and the model actually used at each scale, with the model's name, version, quantisation and
licence, and the measured throughput.}

\section{Computational study on synthetic instances}\label{sec:experiments}

This section measures the algorithms and checks the results of Sections~\ref{sec:structure}
to~\ref{sec:knowledge} on instances drawn at random with fixed seeds. \emph{None of its numbers describes the repair
centre.} Monetary amounts are in arbitrary units. Every table and figure is regenerated by one command
(Appendix~\ref{app:repro}); computing times are those of the machine on which the command was last run.

\subsection{Exactness and cost of pricing}

The implementation solves \eqref{eq:P} by bisection on the labour price, pricing each type by
Remark~\ref{rem:twolevel}. In the test suite its optimal value equals that of an exact linear-programming solver
within $10^{-6}$ on 600 random instances, including instances with missing data, a negative treatment cost and a
rough opening longer than the careful one, and on unit types of nine and ten components where the pricing is done
without enumeration. Table~\ref{tab:pricing} and Figure~\ref{fig:pricing} compare the pricing of one unit type, at
\SynPrixValeurs{} values of the labour price, by enumeration of its plans and by Proposition~\ref{prop:pricing}.
The two agree to $\SynPrixEcart$ on every instance where both run. Enumeration grows by a factor of three per
component and is out of reach beyond a dozen; the separable pricing handles a unit of \SynPrixGrand{} components in
\SynPrixGrandTemps{} milliseconds, most of which is set-up.

\begin{table}[H]
\centering\small
\caption{Pricing one unit type at \SynPrixValeurs{} labour prices: time per type, set-up included (milliseconds).
Synthetic instances.}\label{tab:pricing}
\begin{tabular}{rrrr}
\toprule
Components & Plans & Enumeration & Proposition~\ref{prop:pricing} \\
\midrule
\corps{t_prix}
\bottomrule
\end{tabular}
\end{table}

\begin{figure}[H]
\centering
\includegraphics[width=0.48\textwidth]{\genere/fig_prix.pdf}\hfill
\includegraphics[width=0.48\textwidth]{\genere/fig_variabilite.pdf}
\caption{Left: time to price one unit type, by enumeration and by Proposition~\ref{prop:pricing} (logarithmic
scales; values in Table~\ref{tab:pricing}). Right: variability gap $\Gamma$ as a share of the fluid value, against
the dispersion of weekly arrivals, for three levels of available hours (values in
Table~\ref{tab:gap}).}\label{fig:pricing}
\end{figure}

\subsection{How much can a dynamic policy gain?}

Proposition~\ref{prop:fluid} bounds the gain of any dynamic policy by the variability gap $\Gamma$.
Table~\ref{tab:gap} computes its sample version on a line of \SynGammaTypes{} unit types over
\SynGammaSemaines{} weeks, for arrivals of the same mean and increasing dispersion, and for hours equal to 0.5, 0.8
and 1.2 times the hours the average flow would use without a limit. With a constant flow the gap is zero, as it
must be. It grows with dispersion and falls as hours become plentiful, and it reaches \SynGammaMax{} of the fluid
value in the most dispersed, most constrained case. The reading for a manager is direct: with Poisson-like
arrivals and hours close to the need, a buffer or a forecast has little to recover; with bursty arrivals and
scarce hours, it has much, and the gap says how much before any dynamic policy is built.

\begin{table}[H]
\centering\small
\caption{Variability gap $\Gamma$ as a share of the fluid value $V(\bar A,H)$, by arrival law and available hours
(as a multiple of the unconstrained need at the average flow). Synthetic instance.}\label{tab:gap}
\begin{tabular}{lrrrr}
\toprule
 & Coefficient of variation & \multicolumn{3}{c}{Hours / need} \\
\cmidrule(lr){3-5}
Weekly arrivals & of total arrivals & 0.5 & 0.8 & 1.2 \\
\midrule
\corps{t_variabilite}
\bottomrule
\end{tabular}
\end{table}

\subsection{Learning a yield that the prior gets wrong}

Table~\ref{tab:learning} and Figure~\ref{fig:regret} report a one-component benchmark in which the prior mean of the
yield is \SynTSprior, the true yield \SynTSvrai{} and the reuse threshold of Proposition~\ref{prop:threshold}
\SynTSseuil: reuse is the right decision and the prior says otherwise. The rule that follows the prior and the
greedy rule that follows the posterior mean never test the component, never learn, and lose the same amount on
every copy. Posterior sampling tests almost every copy, recovers the yield and keeps a regret that flattens. The
benchmark is deliberately minimal: it isolates the failure of greedy rules that Section~\ref{sec:learning}
describes. A benchmark on full instances of \eqref{eq:P} is left for future work.

\begin{table}[H]
\centering\small
\caption{One-component learning benchmark (true yield \SynTSvrai, prior mean \SynTSprior, threshold
\SynTSseuil; the oracle earns \SynTSoracle{} per copy). Synthetic.}\label{tab:learning}
\begin{tabular}{lrrr}
\toprule
Policy & Regret per copy & Copies tested & Final estimate of the yield \\
\midrule
\corps{t_apprentissage}
\bottomrule
\end{tabular}
\end{table}

\begin{figure}[H]
\centering
\includegraphics[width=0.48\textwidth]{\genere/fig_regret.pdf}
\caption{Cumulative regret on the benchmark of Table~\ref{tab:learning}. The greedy rule coincides with the rule
that follows the prior: neither ever tests the component.}\label{fig:regret}
\end{figure}

\subsection{Fusion and the reliability of sources}

Table~\ref{tab:fusion} compares the simulated error of the fusion rule with the bound of
Proposition~\ref{prop:fusion}, for \SynFusionClasses{} possible classes. The bound is valid throughout and
conservative by up to an order of magnitude, the price of the union over classes; it is vacuous for one or three
assertions of a mediocre source, which is the regime in which the chain sends a component to review rather than
accept it. For an error bound of 2\%, a source of reliability 0.9 needs \SynFusionKhaut{} independent assertions and
a source of reliability 0.6 needs \SynFusionKbas.

\begin{table}[H]
\centering\small
\caption{Probability that the most probable class is not the true one after $K$ independent assertions of a source
of known reliability $r$ (\SynFusionClasses{} classes; ties count as errors; \SynFusionEssais{} simulated cases per
row), and the bound of Proposition~\ref{prop:fusion}.}\label{tab:fusion}
\begin{tabular}{rrrr}
\toprule
$r$ & $K$ & Simulated error & Bound \eqref{eq:rate} \\
\midrule
\corps{t_fusion}
\bottomrule
\end{tabular}
\end{table}

For Proposition~\ref{prop:strata}, consider a source with \SynStrateNacc{} pieces of evidence on accepted
components, of accuracy \SynStrateRacc, and \SynStrateNfile{} on queued components, of accuracy \SynStrateRfile:
its true accuracy is \SynStrateVraie. An audit validates \SynStrateKacc{} of the former at random, and the
reviewers \SynStrateKfile{} of the latter. Counting all validations together gives \SynStrateEnsemble{} on average
(the proposition predicts \SynStrateAttendu), an error of root mean square \SynStrateEqmEnsemble: the source would
be judged much worse than it is, and components it classifies correctly would be sent to review for nothing. The
stratified estimator gives \SynStrateStrates{} on average, with a root-mean-square error of \SynStrateEqmStrates.

\subsection{Where the acquisition effort goes}

Table~\ref{tab:chain} runs the acquisition chain on a synthetic bill of materials of \SynChainePieces{} components
whose annual masses are log-normal, so that a few components carry most of the stake, with simulated sources of
fixed reliability: a rule on the description, a simulated language model reading descriptions and documents, and a
retrieval service that finds a document for some components and a manufacturer's declaration for a few. Costs are
the working assumptions of the implementation: \CoutDescription{} per reading of a description, \CoutCollecte{}
per retrieval, \CoutExtraction{} per reading of a document, \CoutControle{} per verdict, and \CoutRevue{} per human
review; the acceptance threshold is \SeuilAcceptation.

The chain makes \SynChaineAppels{} model calls, against \SynChaineExhaustif{} for a pipeline that reads, retrieves
and checks every component (\SynChaineEconomie{} fewer), for a total cost of \SynChaineCout{} against a stake
ceiling of \SynChaineEnjeu. No component is charged more than its ceiling: the highest ratio of
spending to ceiling is \SynChaineRatio{} (Lemma~\ref{lem:budget}). It accepts \SynChaineAcceptees{} classes without review, of which
\SynChaineJustesse{} are right, and they carry \SynChaineEnjeuAcceptees{} of the stake. It sends
\SynChaineRevue{} components to a human, which carry \SynChaineEnjeuRevue{} of the stake, and leaves
\SynChaineLaissees{} whose class would earn less than a review costs. The table shows the mechanism: effort follows stake. In the top fifth, \SynChaineHautTraites{} of the
components are settled automatically or sent to a reviewer, and the rest are components whose class, as currently
believed, would earn less than a review costs. In the bottom fifth, \SynChaineBasLaissees{} are left, at an average
cost of \SynChaineBasCout{} per component, and nobody is asked. With these costs the scarce resource is not the
language model, whose calls cost cents, but the human reviewer, and the stake is what rations it.

\begin{table}[H]
\centering\small
\caption{The acquisition chain on a synthetic bill of materials, by fifth of components ranked by stake ceiling
(first fifth: highest). Shares of the fifth's components that are accepted without review, sent to human review,
or left without a class.}\label{tab:chain}
\begin{tabular}{crrrrrrr}
\toprule
 & Share of & Median & Model calls & Cost per & & Sent to & \\
Fifth & total stake & ceiling & per component & component & Accepted & review & Left \\
\midrule
\corps{t_chaine}
\bottomrule
\end{tabular}
\end{table}

\section{Case study: a repair centre for medical equipment}\label{sec:case}

\bandeau

\subsection{Site and data}

The repair centre receives field-replaceable units of medical imaging systems. \apres{one paragraph on the site:
families of equipment, yearly volumes of repaired and scrapped units, the bulk contract and what it costs, the
people who use the system, and since when.} The system reads the centre's operational extractions as they are
exported: repair and scrap records, parts consumption, weighing of the waste bins, a file of early failures of
installed parts, and the bill of materials. Parameters that no extraction contains (operator times, values of
reused parts, recycler prices) are entered in a workbook as ranges with their provenance; estimates are labelled as
such, and a parameter left empty stays unknown. Unit types and components are not named in this paper.
\noteechelle

The line model covers \CasTypes{} unit types, of which \CasIncomplets{} could not be evaluated for lack of
parameters and are reported as such rather than counted as zero. The pipeline ended with \CasAnomalies{} failed
computations. \apres{the verification of the estimated input flow against the weighed mass: number of weeks
compared, share within the stated margin, and what was corrected.}

\subsection{The decision}

Table~\ref{tab:case-decision} gives the indicators on which the decision to install the line rests, at the centre of
the parameter ranges, and the distribution of the annual result when the parameters are drawn from their ranges.
The optimal plan splits \CasPartages{} unit type between two plans, in line with Proposition~\ref{prop:split}.
\apres{which unit types are opened and which are left, which components are reused and which are sorted, and
whether the plan matches what the technicians expected; who validated the valuation of reused parts.}

\begin{table}[H]
\centering\small
\caption{Case study: decision indicators. \ifnum\DonneesDemo=1\relax\textbf{Demonstration data, not
results.}\fi}\label{tab:case-decision}
\begin{tabular}{lr}
\toprule
Indicator & Value \\
\midrule
Variable margin of the line (per year) & \CasMarge \\
Fixed costs of the line, annualised & \CasFixes \\
Annual result against the current contract, $V-K$ & \CasResultat \\
Current cost of bulk treatment (per year) & \CasCoutActuel \\
Operator hours used / available (per year) & \CasHeures{} / \CasCapacite \\
Value of one more operator hour, $\mu^\ast$ & \CasMu \\
Units processed (per year) & \CasUnites \\
Break-even volume (units per year) & \CasSeuil \\
Payback (years) & \CasRetour \\
Mass diverted from the bulk stream / mass scrapped (kg per year) & \CasDetournee{} / \CasMasse{} (\CasPartDetournee) \\
\midrule
Probability that the line is profitable (parameter ranges) & \CasProba \\
Annual result: 5th percentile / median / 95th percentile & \CasPcinq{} / \CasMediane{} / \CasPnonantecinq \\
\bottomrule
\end{tabular}
\end{table}

Table~\ref{tab:case-hours} varies the operator hours. It displays Proposition~\ref{prop:dual} on the site's data:
the margin is concave in the hours, and the value of an additional hour falls by steps until the hours cease to
bind. \apres{the staffing decision this table supported.}

\begin{table}[H]
\centering\small
\caption{Case study: sensitivity to operator hours. \ifnum\DonneesDemo=1\relax\textbf{Demonstration
data.}\fi}\label{tab:case-hours}
\begin{tabular}{rrrrr}
\toprule
Hours per year & Variable margin & Annual result & Occupancy & Value of one more hour \\
\midrule
\corps{t_cas_heures}
\bottomrule
\end{tabular}
\end{table}

\subsection{Is a dynamic policy worth building?}

Replayed week by week on the \CasPeriodes{} weeks of history, with the fixed cost paid every week and no stock, the
line earns \CasMyope{} per year, against \CasFluide{} for the model on the average flow. The variability gap of
Proposition~\ref{prop:fluid} is therefore \CasGamma{} per year, \CasGammaRel{} of the fluid value;
\CasPertes{} weeks end at a loss. \apres{the conclusion drawn: whether the gap justifies a buffer of scrapped
units, and what was decided.}

\subsection{Yields}

The posterior on yields was scored prequentially on \CasLots{} teardown batches (\CasCopies{} copies tested).
Table~\ref{tab:case-yields} compares it with the prior alone and with the raw observed frequency. \apres{which
priors the audit corrected most, in which direction, and whether the three biases named in
Section~\ref{sec:learning} showed.}

\begin{table}[H]
\centering\small
\caption{Case study: prequential scores of the yield estimate (lower is better, except coverage, whose target is
90\%). \ifnum\DonneesDemo=1\relax\textbf{Demonstration data.}\fi}\label{tab:case-yields}
\begin{tabular}{lrrrr}
\toprule
Estimate & Brier score & Log loss & Absolute error & Coverage of the 90\% interval \\
\midrule
\corps{t_cas_rendements}
\bottomrule
\end{tabular}
\end{table}

\subsection{Arrivals}

Table~\ref{tab:case-forecast} compares the forecasting methods of Section~\ref{sec:forecast} by rolling origin on
the weekly series of scrapped units. \apres{the reading of the table; whether the structural model, which uses
the inflow, beats the univariate methods; and the result of any zero-shot foundation model run on the same
series.}

\begin{table}[H]
\centering\small
\caption{Case study: forecasts of weekly scrapped units, rolling origin, all methods scored on the same
forecasts. Root mean squared error relative to the historical mean; pinball loss averaged over quantiles; coverage
of the 80\% interval. \ifnum\DonneesDemo=1\relax\textbf{Demonstration data.}\fi}\label{tab:case-forecast}
\begin{tabular}{lrrrr}
\toprule
Method & Forecasts & Relative RMSE & Pinball loss & Coverage (80\%) \\
\midrule
\corps{t_cas_previsions}
\bottomrule
\end{tabular}
\end{table}

\subsection{Composition of components}

The bill of materials of the modelled units lists \CasCompoPieces{} components. After the run, \CasCompoValidees{}
have a class validated by a human, \CasCompoAcceptees{} a class accepted without review and \CasCompoFile{} wait
in the review queue. The chain made \CasCompoAppels{} calls to the language model in this run, at a cost of
\CasCompoCout; \CasCompoSansAppel{} components were handled without any paid acquisition.
\apres{the model used and where it ran; the measured reliability of each source type after the audit; the number of
reviews done and the time they took; the result of the instruction search on the validated parts (held-out
accuracy before and after, with its test); and one or two examples, anonymised, of components whose class changed
a decision.}

\subsection{What the agents produced}

\apres{the decision note written by the agents on the site's data: how many tool calls the pilot made, how many
numbers the grounding check flagged, the reviewer's reservations, how the note compares with the one a project
engineer wrote, and what the manager did with it.}

\section{Discussion}\label{sec:discussion}

\subsection{Managerial insights}

The insights below follow from the model and do not depend on the numbers of one site.

\emph{One price coordinates the line.} The value of an operator hour is the only number that the bench, the planner
and the manager need to share. The bench uses it through the yield threshold of
Proposition~\ref{prop:threshold}; the planner through the plan of each type; the manager as the worth of an
additional hour of staffing. When hours are plentiful the price is zero and every component is judged on its own
merits; as hours tighten, the same components face a higher bar.

\emph{Measure the gap before building a dynamic policy.} Irregular arrivals invite buffers, forecasts and learned
policies. Proposition~\ref{prop:fluid} gives, from the history alone, the most any of them could add. A site should
compute this number first. If it is small, the myopic weekly plan is enough and the engineering effort is better
spent elsewhere.

\emph{Greedy rules do not learn what they do not try.} A pessimistic prior on a yield is self-confirming under a
rule that follows the mean. Either teardown audits are planned as a deliberate investment, or the decision rule
must randomise in proportion to its uncertainty, as posterior sampling does.

\emph{Spend on knowledge in proportion to what it can change.} Most components of a bill of materials are too light
or too rare for their composition to matter; a few matter a great deal. The stake of
Definition~\ref{def:stake} separates them using only the mass, the volume and the price list, before any model is
called or any person asked. In our experiments the binding resource was the human reviewer, not the language
model: the purpose of the budget is less to save model calls than to decide who deserves a person's attention.

\emph{Treat a language model as a source whose reliability you measure.} The system never accepts a class because a
model said so. It accepts it when the evidence, weighted by reliabilities learned from human validations, makes an
error cheaper than a review. This requires auditing a random sample of what was accepted
(Proposition~\ref{prop:strata}): without the audit, the only validations are of contested components, and every
source looks worse than it is.

\emph{Privacy is an architectural property.} Because outward queries are built before any model call and the model
runs on the owner's hardware, confidentiality does not depend on the model's behaviour or on a provider's terms.
The cost is a smaller model than hosted services offer; the design compensates by asking the model only narrow,
checkable questions and by never trusting the answer.

\subsection{Limitations}

\emph{Evidence.} \ifnum\DonneesDemo=1\relax The case study of this draft runs on fictional data. \fi The
mathematical results are independent of the data; every statement about the site depends on parameters whose
ranges were set by people and on a pilot whose realised margin must be compared with the predicted one.

\emph{Model.} The implementation covers three dispositions and two access levels;
Proposition~\ref{prop:pricing} holds for any finite sets, but the general routine is not coded.
Assumption~\ref{as:indep} is doubtful: a unit scrapped for one cause often has correlated damage, and a model with
unit-level condition classes, in the spirit of quality grading \citep{ferguson2009}, would relax it at the price of
more parameters to learn. Demand for reused parts is assumed to absorb the supply; a cap on the copies returned to
service adds one constraint per saturated component and, by Proposition~\ref{prop:split}, one possible split type
each. The decision maker is risk-neutral; a chance constraint or a conditional value-at-risk on the simulated
distribution of the result would serve one who must not lose money on the line.

\emph{Learning.} We claim no regret guarantee for the posterior-sampling scheme. The dynamic model with inventory
is stated but not solved; the fluid bound is the benchmark against which approximate policies would be measured.
No time-series foundation model was run.

\emph{Knowledge layer.} Independence between sources holds only as far as provenance is traced; reliabilities are
per source type, not per site or per model version; two source types can share a cause of error. The stratified
estimate is unbiased only if the audit is random among all components the chain shows to nobody and the components
sent to review are all reviewed; it is also noisier than a pooled count.
Assumption~\ref{as:loss} ignores contamination penalties. Composition is a single class per component; mass
fractions over several materials call for a Dirichlet belief and a price that is a weighted sum, which leaves
Proposition~\ref{prop:evpi} unchanged. The chain follows a fixed order of actions rather than ranking them by value
per unit cost. Text matching for analogy is a bag of words.

\emph{Agents.} The grounding check guarantees that the numbers of the note exist in the results, not that the
prose around them is right. The instruction search and the agents were developed against stand-in models;
\apres{their behaviour with the model actually deployed, including failure rates of structured output and the
effect of the instruction search.}

\subsection{From the repair centre to production}

The same structure appears in production. A nonconforming assembly at a production stage is a scrapped unit; its
sub-assemblies can be reworked, harvested for new builds, sorted by material, or scrapped as they are.
Table~\ref{tab:production} gives the correspondence. What changes is absorbed by the framework: more dispositions
and deeper bills of materials leave Proposition~\ref{prop:pricing} linear in the number of components; a
component whose class requires a licensed route has a restricted set of dispositions, which removes plans and
leaves \eqref{eq:P} unchanged in form, and is one more reason to know the composition before deciding; with $M$
coupling resources, Proposition~\ref{prop:split} gives at most $M$ split types and the dual has one price per
resource; arrivals are the output of a stage times its nonconformity rate, which is the structural model
\eqref{eq:structural}. None of this has been validated on production data.

\begin{table}[H]
\centering\small
\caption{Correspondence between the repair centre and a production line.}\label{tab:production}
\begin{tabular}{p{0.2\textwidth}p{0.34\textwidth}p{0.36\textwidth}}
\toprule
Abstract object & Repair centre & Production \\
\midrule
Unit type $i$ & type of scrapped unit & nonconforming assembly at a given stage \\
Arrivals $A_{ik}$ & scrapped units per week & output of the stage times its defect rate \\
Inflow $E_{ik}$ & units reaching a repair decision & units produced at the stage \\
Component $j$ & part of the unit & sub-assembly or material fraction \\
Dispositions $\cD$ & reuse, sorted recycling, leave & rework, harvest, sorted recycling, scrap as is \\
Access levels & rough or careful opening & disassembly depth, possibly irreversible \\
Yield $q_{ij}$ & share of extracted parts passing test & share of harvested sub-assemblies that requalify \\
Resource $H$ & operator hours outside the workload & rework stations, test benches \\
Class $C_j$ & recycler price and waste route & recovery route, handling of regulated materials \\
Fixed cost $K$ & teardown and test line & rework cell \\
\bottomrule
\end{tabular}
\end{table}

\subsection{Conclusion}

Disposition of scrapped units is usually presented as an optimisation problem with known data. In the application
that motivated this paper the data were the problem. The system described here couples a disposition model whose
structure makes it cheap to solve and easy to explain, a learning scheme that corrects the priors it starts from,
and a knowledge layer that buys missing facts only when they can pay for themselves and that measures how far each
of its sources can be trusted. A language model has a place in such a system: hosted on the owner's machine,
confined to narrow questions and a closed list of tools, treated as a fallible source, and checked by programs
where it writes numbers and by people where it matters. The pieces are simple. What we believe is of general use is
the way they are tied together by one quantity, the value that a fact or an hour has for the decision at hand.

\paragraph{Disclosure.} The software, the proofs and this draft were prepared with the assistance of a generative
language-model system, under the direction of the authors. The propositions and their proofs were checked by
independent automated reviews and by numerical tests, which is not peer review; the authors are responsible for
verifying every statement, proof and reference before submission. \apres{acknowledgements, and the statement on
data availability agreed with the organisation.}


\singlespacing
\bibliographystyle{plainnat}
\bibliography{references}

\appendix
\section{Proofs}\label{app:proofs}

\subsection*{Proposition~\ref{prop:pricing}}
Write $G_i(\ell)$ for the term in braces in \eqref{eq:pricing}. ($\le$) Let $s$ be any plan and $\ell=\ell(s)$.
Every $s_j$ satisfies $\ell_{s_j}\le\ell$, so $n_{ij}(\rho_{ijs_j}-\omega_it_{js_j})$ is at most the inner maximum,
and $g_{is}-\mu h_{is}=-\omega_ia_{i\ell}+\sum_jn_{ij}(\rho_{ijs_j}-\omega_it_{js_j})\le G_i(\ell)$. ($\ge$) Fix
$\ell$ and let $s^\ell$ pick, for each $j$, a maximiser of the inner problem. Then $\ell(s^\ell)\le\ell$, hence
$a_{i\,\ell(s^\ell)}\le a_{i\ell}$ by Assumption~\ref{as:nested}, and since $\omega_i\ge0$,
$g_{is^\ell}-\mu h_{is^\ell}\ge G_i(\ell)$. The count of operations is that of the double maximum. \qed

\subsection*{Remark~\ref{rem:twolevel}}
Partition the plans by induced access level. Level 0 contains only $s^0$, of reduced margin 0. A plan of level 1
sorts a nonempty set of components and reuses none; its reduced margin is $-\omega_ia_{i1}$ plus the sum of
$\Psi^{\Sd}_j$ over the sorted components, maximised by sorting those with $\Psi^{\Sd}_j\ge0$; if there is none,
every plan of this level is worth less than 0 and the term can be dropped. A plan of level 2 reuses a nonempty set.
Ignoring that requirement, each component takes $\max(\Psi^{\R}_j,\Psi^{\Sd}_j,0)$; if at least one component
attains this maximum with reuse the requirement holds at no cost, and otherwise the cheapest way to satisfy it is
to switch to reuse the component for which the loss $\max(\Psi^{\Sd}_j,0)-\Psi^{\R}_j>0$ is least. The maximum of
the three cases is $\varphi_i(\mu)$; the opening times enter only through the level actually induced, so no
ordering between them is used. \qed

\subsection*{Proposition~\ref{prop:dual}}
Dualise the capacity constraint with multiplier $\mu\ge0$. The remaining problem separates by type and its value is
$N_i\varphi_i(\mu)$, since all of $N_i$ goes to a best plan. Strong duality holds because \eqref{eq:P} is a
feasible bounded linear programme. The optimal value of a maximisation linear programme is concave and piecewise
linear in its right-hand side; positive homogeneity follows from scaling $x$. Monotonicity in $H$ is the
relaxation of a constraint; monotonicity in $N_i$ holds because additional units can be assigned $s^0$ at no cost.
For any $H'$, weak duality gives
$V(N,H')\le\mu^\ast H'+\sum_iN_i\varphi_i(\mu^\ast)=V(N,H)+\mu^\ast(H'-H)$; applied to $H'>H$ and $H'<H$ this
places $\mu^\ast$ between the two one-sided derivatives of $V$ in $H$. If the budget is not exhausted at an
optimum, complementary slackness gives $\mu^\ast=0$. \qed

\subsection*{Proposition~\ref{prop:split} and Corollary~\ref{cor:frontier}}
Types with $N_i=0$ have all their variables equal to zero; delete them, and let $\cI^+=\{i:N_i>0\}$. Add a slack to
each of the $M$ coupling constraints. The programme has $|\cI^+|+M$ equality constraints, so an optimal basic
solution, which exists because the programme is feasible and bounded, has at most $|\cI^+|+M$ positive variables.
Each type of $\cI^+$ needs at least one positive plan variable, which leaves at most $M$ further positive
variables, hence at most $M$ types with a second plan. For the second claim, with $M=1$, let $(\zeta,\mu^\ast)$ be
optimal dual variables: dual feasibility reads $\zeta_i\ge g_{is}-\mu^\ast h_{is}$ for all $s$, and complementary
slackness gives equality whenever $x_{is}>0$. A plan in use therefore maximises $g_{is}-\mu^\ast h_{is}$ over
$\cS_i$: the line of slope $\mu^\ast$ through $(h_{is},g_{is})$ supports the point set from above, and two plans of
the same type in use lie on the same supporting line. For the corollary, a weighted linear outcome only changes
the coefficients $g$; $M$ lower bounds on linear outcomes are $M$ further coupling constraints, $M+1$ with the
labour budget. \qed

\subsection*{Remark~\ref{rem:integrality}}
Let the split type use plans $s$ and $s'$ with $h_{is}\le h_{is'}$ in quantities that sum to the integer $N_i$.
Rounding the quantity of $s'$ down and that of $s$ up by the same fraction does not increase the hours and changes
the objective by less than $|g_{is'}-g_{is}|$. \qed

\subsection*{Proposition~\ref{prop:threshold}}
Per copy and with indices dropped, the reduced values are
$\psi^{\R}=q\bar v+(1-q)mp+\tau m-\kappa-\omega t_{\R}$, $\psi^{\Sd}=m(p+\tau)-\omega t_{\Sd}$ and $\psi^0=0$.
$\psi^{\R}\ge\psi^{\Sd}$ rearranges to $q(\bar v-mp)\ge\kappa+\omega(t_{\R}-t_{\Sd})$, and $\psi^{\R}\ge0$ to
$q(\bar v-mp)\ge\kappa+\omega t_{\R}-m(p+\tau)$. Both hold if and only if $q(\bar v-mp)$ is at least the larger
right-hand side, which is $\kappa+\omega t_{\R}-\min\{\omega t_{\Sd},m(p+\tau)\}$; divide by $\bar v-mp>0$. The
slope of the numerator in $\omega$ is $t_{\R}-t_{\Sd}\ge0$ where the minimum is $\omega t_{\Sd}$, and
$t_{\R}\ge0$ elsewhere. \qed

\subsection*{Proposition~\ref{prop:fluid}}
Fix $T$ and let $\bar x=\frac1T\sum_{k\le T}x_k$. With an empty initial stock, units processed by $T$ were received
by $T$, so $\sum_s\bar x_{is}\le\frac1T\sum_{k\le T}A_{ik}$ for every $i$; each period uses at most $H$ hours, so
$\sum_{i,s}h_{is}\bar x_{is}\le H$. This holds whatever the policy knows, including future arrivals. Completing
$\bar x$ with the plan $s^0$, which costs no hours and earns nothing, gives a feasible point of \eqref{eq:P} with
data $(\frac1T\sum_{k\le T}A_k,H)$ and the same objective: this is the pathwise claim. Taking expectations and
using the concavity of $V$ in $N$ (Proposition~\ref{prop:dual}) with Jensen's inequality gives
$\frac1T\E\sum_kg^\top x_k\le V(\bar A^T,H)$ with $\bar A^T=\frac1T\sum_{k\le T}\E A_k$, and \eqref{eq:fluid}
follows by continuity of $V$, which is finite, concave and piecewise linear on $N\ge0$. The claim on the myopic
policy is Jensen's inequality for $N\mapsto V(N,H)$, applied to the empirical distribution over periods and to the
expectation. \qed

\subsection*{Proposition~\ref{prop:evpi}}
For a fixed feasible $x$ the objective $g(\theta)^\top x$ is affine in $\theta$, and $V$ is the pointwise maximum of
these affine functions over a feasible set that does not depend on $\theta$, hence convex. Let $x^\ast$ be optimal
at $\bar\theta$. Then $V(\bar\theta)=g(\bar\theta)^\top x^\ast=\E_F[g(\theta)^\top x^\ast]\le\E_F[V(\theta)]$, with
equality if and only if $g(\theta)^\top x^\ast=V(\theta)$ almost surely. For the bound, let the slopes of $V$ lie
in $[\underline\lambda,\overline\lambda]$ with $\overline\lambda-\underline\lambda=\Delta$, and put
$\lambda^\circ=(\underline\lambda+\overline\lambda)/2$. The function $U(\theta)=V(\theta)-\lambda^\circ\theta$ is
convex with slopes in $[-\Delta/2,\Delta/2]$, hence $|U(\theta)-U(\bar\theta)|\le\tfrac12\Delta|\theta-\bar\theta|$;
and $\E_F[V(\theta)]-V(\bar\theta)=\E_F[U(\theta)-U(\bar\theta)]$ because $\E_F[\theta-\bar\theta]=0$. The constant
is attained by $V(\theta)=\Delta\pos{\theta-\bar\theta}$, for which $\EVPI=\Delta\,\E_F\pos{\theta-\bar\theta}=\tfrac12\Delta\,\E_F|\theta-\bar\theta|$. For a price,
the slope of the margin of any feasible $x$ is $m_j$ times the number of copies of $j$ sold by the kilogram, which
lies in $[0,N^{c}_jm_j]$ since $j$ may belong to several types. For a yield, it is $\lambda_{ij}$ times the number
of copies of $j$ reused in type $i$, which lies between $0$ and $N_in_{ij}\lambda_{ij}$. \qed

\subsection*{Proposition~\ref{prop:class}}
Fix a price $p$ and let $x$ be optimal for it. Replace, in every plan, the disposition $\Sd$ of $j$ by $0$. The
access level of a plan does not rise: it stays 2 if another component is reused, stays 1 if another is sorted, and
otherwise falls to 0, where the opening time is zero. Hours do not increase, so the new solution is feasible for
the problem without a price. Per unit of type $i$, a plan that sorted $j$ loses $n_{ij}m_j(p+\tau)$ minus the
labour saved; a plan that reused $j$ loses $n_{ij}(1-q_{ij})m_j(p+\tau)$, the difference between the two
expressions of $\rho_{ij\R}$ in Section~\ref{sec:model}; other plans lose nothing. Each loss is at most
$n_{ij}m_j\pos{p+\tau}$, and at most $N_i$ units of type $i$ are processed, so
$V(p)-V(\out_j)\le N^{c}_jm_j\pos{p+\tau}$. Take the expectation over the class; for the value $\out$ there is no
sorted outlet and both sides are zero. Conversely, when $p+\tau\ge0$ every plan is worth at least as much with the
price as without, so $V(p)\ge V(\out_j)$. \qed

\subsection*{Lemma~\ref{lem:budget}}
By induction on the paid actions. Initially $S_j=0\le\bar e_j$. A paid action $a$ is performed only if the stake is
known and $e_j(\pi)\ge S_j+\phi_a$ for the belief $\pi$ held at that moment; after it, the new cumulative cost is
$S_j+\phi_a\le e_j(\pi)\le\bar e_j$, because $\pi$ is a probability vector and $v\ge0$. An action that fails for a
technical reason is not charged, and $S_j$ is read from the journal, so the invariant holds across runs. If
$\bar e_j<\min_a\phi_a$ or the stake is unknown, the test fails for every action. \qed

\subsection*{Proposition~\ref{prop:screen}}
Under Assumption~\ref{as:noerror}, whatever the procedure decides, the realised value is at most
$\max\{V(p(C_j)),V(\out_j)\}$, so $G\le\pos{V(p(C_j))-V(\out_j)}\le M_jv(C_j)$, by the first inequality of
Proposition~\ref{prop:class} and its extension to the other streams (the right-hand side is nonnegative, and both
sides vanish for the value $\out$). Taking the expectation under $\pi$ gives $\E_\pi[G]\le e_j(\pi)$. (i) and (ii)
follow: if $\E_\pi[G]-\E_\pi[\Phi]>0$ then $\E_\pi[\Phi]<e_j(\pi)$, and for a single action $\Phi=\phi_a$. (iii)
The belief after an observation is the posterior of the model of Section~\ref{sec:belief}; under that model's
predictive distribution the expected posterior equals the prior, $\E[\pi_{t+1}(c)\mid\pi_t]=\pi_t(c)$, and $e_j$ is
linear in $\pi$. \qed

\subsection*{Lemma~\ref{lem:accept}}
If the true class is $\hat c$ nothing is lost. If it is $c\ne\hat c$, the loss is at most $M_jv(c)$ by
Assumption~\ref{as:loss}. The expected loss is therefore at most
$M_j\sum_{c\ne\hat c}\pi(c)v(c)\le(1-\pi(\hat c))M_j\max_cv(c)=R_j$. A review reveals the class at cost
$\phi_{\mathrm{rev}}$ and loses nothing else, so accepting is cheaper in expectation whenever
$R_j<\phi_{\mathrm{rev}}$. Under calibration the probability that an accepted class is wrong is
$1-\pi(\hat c)\le1-\theta$. The implementation restricts the maximum in $R_j$ to classes of probability at least
$0.01$, which changes the bound by at most $0.01\,n\,\bar e_j$. \qed

\subsection*{Proposition~\ref{prop:fusion}}
If $C_j=c$, the probability of the $K$ assertions is $r^{k_c}q^{K-k_c}$, which with a uniform prior gives a
posterior proportional to $(r/q)^{k_c}$; the three cases follow from the sign of $\ln(r/q)$, since $r>q$ if and
only if $r>1/n$. Let $c^\ast$ be the true class and $c'\ne c^\ast$. Then $D=k_{c^\ast}-k_{c'}=\sum_{i\le K}X_i$
with independent $X_i$ equal to $+1$, $-1$, $0$ with probabilities $r$, $q$, $1-r-q$. If $r>q$, for every
$\lambda\ge0$,
\[
\Pr(D\le0)\le\E\big[e^{-\lambda D}\big]=\big(re^{-\lambda}+qe^{\lambda}+1-r-q\big)^K ,
\]
and the choice $e^{\lambda}=\sqrt{r/q}\ge1$ gives $\big(1-r-q+2\sqrt{rq}\big)^K=\varrho^K$ (when $q=0$ the event
has probability zero for $K\ge1$, and $\varrho=0$). If $r<q$, the same computation with $-D$ bounds $\Pr(D\ge0)$ by
$\varrho^K$. The most probable class differs from $c^\ast$, or ties with another, only if $D\le0$ (respectively
$D\ge0$) for some $c'$; the union over the $n-1$ classes gives \eqref{eq:rate}, and solving
$(n-1)\varrho^K\le\delta$ gives the number of assertions. Since $\sum_K(n-1)\varrho^K<\infty$, by the
Borel--Cantelli lemma $\hat c_K=c^\ast$ for all large $K$ almost surely; and by the strong law of large numbers
$|D|\to\infty$ with the sign of $r-q$ for each $c'$, so the posterior ratio $(r/q)^{D}\to\infty$ and the posterior
probability of $c^\ast$ tends to one. \qed

\subsection*{Proposition~\ref{prop:strata}}
Under simple random sampling without replacement of $k_s$ pieces among the $n_s$ of stratum $s$, of which
$n_sr_s$ are right, $y_s$ is hypergeometric with mean $k_sr_s$. (i) $\E[\hat r]=\sum_sn_sr_s/n=r$. (ii)
$\E[(y_{\mathsf A}+y_{\mathsf Q})/k]-r=r_{\mathsf A}(k_{\mathsf A}/k-n_{\mathsf A}/n)+r_{\mathsf Q}(k_{\mathsf Q}/k-n_{\mathsf Q}/n)$,
and $k_{\mathsf Q}/k-n_{\mathsf Q}/n=-(k_{\mathsf A}/k-n_{\mathsf A}/n)$. The sign follows: the queue is
over-represented when $k_{\mathsf A}/k<n_{\mathsf A}/n$. (iii) The pieces of stratum $s$ not validated number
$n_s-k_s$, of which $n_sr_s-y_s$ are right; their expected accuracy is $(n_sr_s-k_sr_s)/(n_s-k_s)=r_s=\E[y_s/k_s]$.
Weighting by $n_s-k_s$ on both sides gives the claim. \qed

\subsection*{Lemma~\ref{lem:grounding}}
(i) The pilot's answer is parsed as one object naming a tool; a name outside the list returns an error message
without executing anything, arguments outside those declared for the tool are dropped, and the tools are
functions of the system with no access to a shell or to the network other than through the acquisition chain of
Section~\ref{sec:knowledge}. (ii) After the last model call, a program extracts the numerals of the note and of
the evidence and computes the set of numerals of the note, other than integers below ten, that match no numeral of
the evidence; the warning at the head of the note lists exactly this set, and is absent if and only if the set is
empty. (iii) Unavailability of the model, an unreadable answer twice in a row, or exhaustion of the turns ends the
pilot with no finding, in which case the fixed plan runs the remaining tools and the template note is used; the
same happens if the writer fails. \qed

\subsection*{Lemma~\ref{lem:heldout}}
The returned instruction is a function of the batch and selection sets and of the models' answers on them. The test
parts are independent of these sets and identically distributed, so, given what the search saw, the paired
differences of scores on the $n_T$ test parts are independent, lie in $[-1,1]$ and have mean $\Delta$. The
one-sided inequality of \citet{hoeffding1963} gives
$\Pr(\hat\Delta-\Delta\ge\epsilon)\le\exp(-2n_T\epsilon^2/2^2)$. If $\Delta<-\epsilon$, adoption requires
$\hat\Delta\ge0$, hence $\hat\Delta-\Delta>\epsilon$. The search loop starts an iteration only if the calls made
plus those the iteration can make stay within $B$, each (instruction, example) pair is evaluated once, and the
test evaluates two instructions on $n_T$ parts. \qed

\section{Implementation and reproducibility}\label{app:repro}

The system is one Python package with four dependencies (numerical arrays, data frames, spreadsheet reading,
plotting). It runs offline. The language model is optional and is reached through a local chat-completions
endpoint. Table~\ref{tab:map} maps the paper to the code and to the tests that check it.

\paragraph{Commands.} \texttt{python -m erc.diagnostic} checks the installation without any network access.
\texttt{python -m erc.agent} runs the system end to end and writes the decision note, the evidence and the journal
of the agents; with \texttt{--sans-modele} it runs the fixed plan. \texttt{python -m erc.papier} recomputes every
number, table body and figure of this paper from the data of the installation and writes them to a directory of
the paper's folder, from which this document reads them; with \texttt{--compiler} it also compiles the document.
Nothing in that directory is typed by hand. The synthetic experiments use fixed seeds; computing times
depend on the machine. The case-study values are those of the data the command is run on
\ifnum\DonneesDemo=1\relax (in this draft, the fictional demonstration data)\fi.

\paragraph{What is checked numerically.} The optimum against an exact linear-programming solver on random
instances; the separable pricing against enumeration; Propositions~\ref{prop:dual} to~\ref{prop:class} on random
instances, including the tight case of \eqref{eq:evpi}; the fusion formula on hand-computed cases; the bound of
Proposition~\ref{prop:fusion} against simulation; the bias formula of Proposition~\ref{prop:strata}; the budget
invariant of Lemma~\ref{lem:budget} on simulated bills of materials; the grounding check of
Lemma~\ref{lem:grounding} against a stand-in model that invents numbers; the absence of any internal identifier or
description in outgoing queries; and the absence of leakage from a batch to its own prediction in the prequential
scores.

\begin{table}[H]
\centering\footnotesize
\caption{Map from the paper to the implementation.}\label{tab:map}
\begin{tabular}{p{0.36\textwidth}p{0.24\textwidth}p{0.32\textwidth}}
\toprule
Paper & Module & Tests \\
\midrule
Values, \eqref{eq:P}, multiplier, structure (\S\ref{sec:model}--\ref{sec:structure}) & \texttt{erc/ligne.py} & \texttt{tests/test\_ligne.py} \\
Same model in the browser tool & \texttt{erc/outil\_modele.js} & \texttt{tests/test\_outil.py} \\
Prior, posterior, TS-LP, scores (\S\ref{sec:learning}) & \texttt{erc/rex.py} & \texttt{tests/test\_stats.py} \\
Early failures & \texttt{erc/fiabilite.py} & \texttt{tests/test\_stats.py} \\
Period-by-period replay, gap $\Gamma$ & \texttt{erc/ligne.py} & \texttt{tests/test\_information.py} \\
Forecasting and its evaluation & \texttt{erc/prevision.py} & \texttt{tests/test\_prevision.py} \\
Belief, reliabilities by strata (\S\ref{sec:belief}) & \texttt{erc/cortex.py} & \texttt{tests/test\_composition.py} \\
Acquisition chain, budget, acceptance & \texttt{erc/composition.py} & \texttt{tests/test\_composition.py} \\
Bounds and exact values of information & \texttt{erc/information.py} & \texttt{tests/test\_information.py} \\
Human review and audit sample & \texttt{erc/composition\_outil.py} & \texttt{tests/test\_composition\_outil.py} \\
Self-hosted model, trust policy & \texttt{erc/modele.py} & \texttt{tests/test\_modele.py} \\
Instruction search & \texttt{erc/gepa.py} & \texttt{tests/test\_gepa.py} \\
Agents, containment, grounding (\S\ref{sec:system}) & \texttt{erc/agent.py} & \texttt{tests/test\_agent.py} \\
Offline diagnostic & \texttt{erc/diagnostic.py} & \texttt{tests/test\_diagnostic.py} \\
Numbers of this paper & \texttt{erc/papier.py} & \texttt{tests/test\_papier.py} \\
\bottomrule
\end{tabular}
\end{table}


\end{document}
